\chapter{Aru}\label{ch:Aru}
In this chapter we discuss the position
and internal subgrouping of the \tsc{Aru} languages.
For this, we rely almost entirely on \citet{Nivens2017},
a database of over 900 Proto-\ili{Aru} reconstructions with supporting reflexes.
All data and Proto-\ili{Aru} reconstructions in this chapter are from
this database with some differences in the analysis of the Proto-\ili{Aru} phoneme inventory.
Any future citation of this chapter or its conclusions must also
include a reference to \citet{Nivens2017}.\footnote{
		We suggest
		the following citation format: “Grimes
		and Edwards, based on \citet{Nivens2017}, {\ldots}”}
\citet{Nivens2017}
compiles data from several sources including his own fieldwork
and fieldwork by other linguists, mainly: Jock Hughes (many languages), Benjamin T.
Daigle (\ili{Batuley}, \ili{Mariri}), and Antoinette Schapper (\ili{Ujir}).\footnote{
		None of the people mentioned here are responsible for any
		errors in the presentation of this data in this chapter.
		Neither do they necessarily agree with our interpretation of
		this data.} 
Another important source in understanding Proto-\ili{Aru} is \citet{Nivens1993},
a reconstruction of Proto-\ili{Aru} based on data from West \ili{Tarangan}, \ili{Dobel}, and \ili{Kola}.
\citet[127--133]{Collins1982c} also provides a classification of
the \tsc{Aru} languages and describes the sound changes in the villages from which he collected data.
\citet{Chlenov1976} also presents a classification of the \tsc{Aru}
languages, with a more detailed discussion
and presentation given in \citet{ChlenovChlenova2008a}.

A map of the \tsc{Aru} languages is given in \frf{06-01}. For
languages spoken in more than one village,
the village from which data in this chapter is usually drawn is also indicated with a red dot.
All data which is not from these villages is explicitly flagged
in the table notes.
\ili{Karey}, \ili{Lorang}, \ili{Mariri}, and \ili{Kompane} are only spoken in one village each.

\begin{figure}\centering
	\includegraphics[height=0.96\textheight]{Aru}
	\fcp{\tsc{Aru} languages}{06-01}
\end{figure}

Transcriptions of data are adapted to IPA standards,
with the following exceptions.
The letter ‹y› is used for the palatal glide [j],
\tsc{Tarangan} mid-close [e] and [o] are transcribed ‹ê›
and ‹ô› (mid-open [ɛ]
and [ɔ] are transcribed ‹ɛ›
and ‹ɔ›), and unexpected stress is marked by an acute accent.
(The normal pattern is penultimate stress,
or final stress if the final syllable contains a mid-vowel.)
The symbol ‹ʍ› represents a voiceless labio-velar fricative.

\citet{GrayGreenhill2009} subgroup the \tsc{Aru} languages together with the languages
of “\ili{Central} Maluku” based on a Bayesian phylogenetic analysis
of the lexicon, albeit with marginal support.
Looking at both the lexicons
and historical phonologies of our \tsc{Sula-Buru},
\tsc{\ili{Ambon}-Seram}, and \tsc{Seram-Tanimbar-Bomberai} subgroups,
we cannot support such a link.
Furthermore, while the languages in these three subgroups of
geographical central Maluku overall are fairly conservative in
retaining stable and widespread PMP vocabulary
and phonotactic structure, the \tsc{Aru} languages are not.

Regarding the lexicons, \citet[245]{Blust1993} observes,
“Almost all CMP languages apart from those of the \ili{Aru} Islands
are lexically quite conservative.” In the \tsc{Aru} languages,
several basic numerals, some common body parts,
and common verbs often do not reflect the PMP words widespread in nearby regions.
A sample of replacement innovations in the basic vocabulary in
\tsc{Aru} languages is given in \trf{06-01} to illustrate their lexical divergence.
The source of such lexical replacements is currently unknown.

\begin{table}\tcp{Sample of \tsc{Aru} basic vocabulary replacement}{06-01}
\begin{adjustbox}{width=\textwidth}
	\begin{threeparttable}\st{2pt}
	\begin{tabular}{llllllll}\lsptoprule
(PMP)	&	(*təlu)	&	(*əpat)	&	(*ənəm)	&	(*laŋit)	&	(*qaqay)	&	(*daRaq)	&	(*matay)\\  \midrule
**gloss	&	three	&	four	&	six	&	sky	&	foot/leg	&	blood	&	die\\
Proto-\ili{Aru}	&	**lati	&	**ka(ɸa)\su{‡}	&	**ndumbu	&	**ŋgamáR	&	**amba-	&	**undu	&	**-woy\\ \midrule
\ili{Ujir}	&	{\hr\hr}lati	&	{\hr\hr}ka	&	{\hr\hr}dubu	&		&	{\hr\hr}ay	&	{\hr\hr}udu	&	{\hr\hr}-woy\\
\ili{Kompane}	&	{\hr\hr}lasi	&	{\hr\hr}kaha	&	{\hr\hr}dubu	&	{\hr\hr}ŋamá	&	{\hr\hr}aba-y	&	{\hr\hr}gudu	&	{\hr\hr}-goy\\
East \ili{Kola}	&	{\hr\hr}las	&	{\hr\hr}kafa	&	{\hr\hr}dum	&	{\hr\hr}gamáh	&	{\hr\hr}ay tul	&	{\hr\hr}un	&	\hp{**-}\tcg{(sowih)}\\
\ili{Lola}	&	{\hr\hr}lasi	&	{\hr\hr}ʔaha	&	{\hr\hr}dubu	&	{\hr\hr}ŋamá	&	{\hr\hr}yabi	&	{\hr\hr}wudu	&	{\hr\hr}-woy\\
\ili{Lorang}	&	{\hr\hr}lay\su{†}	&	{\hr\hr}kawa	&	{\hr\hr}dubu	&	{\hr\hr}maŋár\su{Δ}	&	{\hr\hr}yabay	&		&	{\hr\hr}-goy\\
\ili{Dobel}	&	{\hr\hr}lay	&	{\hr\hr}ʔawa	&	{\hr\hr}dubu	&	{\hr\hr}maŋár\su{Δ}	&	{\hr\hr}yabay	&	{\hr\hr}kudu	&	{\hr\hr}-koy\\
\ili{Koba}	&	{\hr\hr}lay\su{†}	&	{\hr\hr}kawa-y	&	{\hr\hr}dubi	&	{\hr\hr}maŋár\su{Δ}	&	{\hr\hr}yabay	&		&	{\hr\hr}-ʍoy\\
\ili{Barakai}	&	{\hr\hr}laʔ	&	{\hr\hr}kaw	&	{\hr\hr}dom	&	{\hr\hr}gamár	&	{\hr\hr}aba	&	{\hr\hr}gon	&	{\hr\hr}-gwiʔ\\
\ili{Karey}	&	{\hr\hr}lado	&	{\hr\hr}kaw	&	{\hr\hr}dom	&	{\hr\hr}ŋamár	&	{\hr\hr}aba	&	{\hr\hr}gun	&	{\hr\hr}-go\\
\ili{Batuley}	&	{\hr\hr}laes\su{†}	&	{\hr\hr}kaw	&	{\hr\hr}dum	&	{\hr\hr}ŋamár	&	{\hr\hr}abey	&	{\hr\hr}gun	&	{\hr\hr}-woy\\
\ili{Mariri}	&	{\hr\hr}laʤ	&	{\hr\hr}kaw	&	{\hr\hr}dom	&	{\hr\hr}ŋamár	&	{\hr\hr}aba-y	&	{\hr\hr}gun	&	{\hr\hr}-wúy\\
\ili{Manombai}	&	{\hr\hr}lasi	&	{\hr\hr}ka	&	{\hr\hr}dubu	&	{\hr\hr}maŋár\su{Δ}	&	{\hr\hr}ebay	&	{\hr\hr}gudu	&	\hp{**-}{(seuri)}\\
NE. \ili{Tarangan}	&	{\hr\hr}lat	&	{\hr\hr}kaw	&	{\hr\hr}dɔm	&	{\hr\hr}maŋár\su{Δ}	&	{\hr\hr}ɛbɛ	&		&	{\hr\hr}-gɔy\\
SE. \ili{Tarangan}	&	{\hr\hr}lat	&	{\hr\hr}kaw	&	{\hr\hr}dɔm	&	{\hr\hr}maŋár\su{Δ}	&		&		&	{\hr\hr}-gɔy\\
NW. \ili{Tarangan}	&	{\hr\hr}lat	&	{\hr\hr}ka	&	{\hr\hr}dum	&	{\hr\hr}maŋár\su{Δ}	&	{\hr\hr}aba-y	&	{\hr\hr}gun	&	{\hr\hr}-gɔy\\
SW. \ili{Tarangan}	&	{\hr\hr}lat	&	{\hr\hr}ka	&	{\hr\hr}dum	&	{\hr\hr}maŋár\su{Δ}	&	{\hr\hr}aba-y	&	{\hr\hr}gun	&	{\hr\hr}-gɔy\\
		\lspbottomrule
	\end{tabular}
		\begin{tablenotes}\tnsize
			\item[†]	{\ili{Lorang} \it{lay},
								\ili{Koba} \it{lay}, and \ili{Batuley} \it{laes} are \tsc{3pl} forms.
								Other forms are inanimate.}
			\item[‡]	{Final (optional) **ɸa of Proto-\ili{Aru} **ka(ɸa) ‘four’ may be connected
								with PMP *əpat ‘four’.}
			\item[Δ]	{Some languages with **ŋg > \it{ŋ} have metathesis of the
								initial and medial consonants of **ŋgamáR ‘sky’ > **ŋamár > **maŋár.}
		\end{tablenotes}
	\end{threeparttable}
\end{adjustbox}
\end{table}

In addition to lexical replacements,
the \tsc{Aru} languages often do not follow the lexical semantic
patterns which are widespread in the region.
Thus, for instance, many languages of the region have shifted
PMP *duRi ‘thorn’ > ‘bone’,
while \tsc{Aru} languages are conservative in having reflexes
of PMP *tuqəlan > **tul- for ‘bone’,
while Proto-\ili{Aru} **ruRi retains the meaning ‘thorn’.

\ili{Aru} languages are also different from the languages of central
Maluku in their historical phonologies.
\tsc{\ili{Ambon}-Seram} languages tend to merge *d/*z/*j -- but not fully.
\tsc{Aru} languages keep the three distinct, but merge *j/*R.
\tsc{\ili{Ambon}-Seram} languages tend to merge *n/*ñ/*ŋ -- but not fully.
\tsc{Aru} languages keep them distinct.
\tsc{\ili{Ambon}-Seram} languages tend to merge *R/*l -- but not fully.
\tsc{Aru} languages keep them distinct.

The \tsc{Aru} languages have undergone many sound changes involving
both conditioned and unconditioned splits, with original conditioning environments
often erased by subsequent changes.
The complexity of the data is well exemplified by \citet{Blust2014}
who examines the historical phonology of \ili{Dobel} (based on the
data in \citealp{Hughes1995}).
Recurrent glide insertion, glide fortition,
and unstressed vowel loss have combined to obscure many regular Austronesian inheritances.
Thus, for instance, \ili{Dobel} \it{say} ‘fire’ is regular from PMP
*hapuy ‘fire’ via several intermediate stages: *hapuy > **api
> **ai > **yai > **ʤai > **sai > \it{say}.

\scp{Proto-Aru}{06-01}
\citet{Nivens1993,Nivens2017} posits the following regular sound changes
between PMP and Proto-Aru: *p > **ɸ,
*b > **p, *d > **r,
*j/*R > **R, *z/*y > **y,
and *q/*h > {\0}, with other consonants usually remaining unchanged.
Proto-Aru **ɸ was either a bilabial fricative [ɸ] or labiodental fricative [f].

In addition to the consonants which are reflexes of PMP proto-phonemes,
\citet{Nivens2017} also reconstructs three voiced stops for which
we cannot identify a consistent
and regular Austronesian source.
Our interpretation of the Proto-Aru phoneme inventory differs
slightly from that proposed by \citet{Nivens1993,Nivens2017}, in that we take the Proto-Aru voiced stops to be prenasalised.
The possible values of the Proto-Aru consonants which have changed
from PMP, as well as those which have no consistent PMP source are given
in \trf{06-02} along with their reflexes in selected languages.
The third line gives the phonetic values assigned to them by
\citet{Nivens1993,Nivens2017}, while the fourth line gives the interpretation which we
take in this chapter.

\begin{table}\tcp{Selected Proto-Aru consonants with selected reflexes}{06-02}
%\begin{adjustbox}{width=\textwidth}
	\begin{threeparttable} \st{4pt}
	\begin{tabular}{lllllllllll}\lsptoprule
PMP	&	\hr*p	&	\hr*b	&	\hr*d	&	*j/*R	&	*z/*y	&	\mc{2}{l}{(*mb)}			&		&		&	(*g)	\\	
Proto-Aru	&	**ɸ	&	**p	&	**r	&	**R	&	**y	&	\mc{2}{l}{**mb\su{‡}}			&	\mc{2}{l}{**nd}			&	**ŋg	\\	
Nivens	&	\hr[ɸ]/\hr[f]	&	\hr[p]	&	\hr[r]	&	\hr[ʀ]\su{†}	&	\hr[y]	&	\mc{2}{l}{\hr[b]}			&	\mc{2}{l}{\hr[d]}			&	\hr[g]	\\	
this chapter	&	\hr[ɸ]	&	\hr[p]	&	\hr[r]	&	\hr[ʁ]	&	\hr[y]	&	\mc{2}{l}{\hr[mb]}			&	\mc{2}{l}{\hr[nd]}			&	\hr[ŋg]	\\	
	&		&		&		&		&		&		&	/{\gap\#}	&		&	/{\gap\#}	&		\\	\midrule
Kompane	&	{\hr\hr}h	&	{\hr\hr}ɸ	&	{\hr\hr}r	&	{\hr\hr}h	&	{\hr\hr}ʤ	&	{\hr\hr}b 	&	-m	&	{\hr\hr}d	&	-n	&	{\hr\hr}ŋ	\\	
\tsc{Kola}	&	{\hr\hr}f	&	{\hr\hr}ɸʷ	&	{\hr\hr}r	&	{\hr\hr}h	&	{\hr\hr}y	&	{\hr\hr}b	&	-m	&	{\hr\hr}d	&	-n	&	{\hr\hr}g	\\	
Lorang	&	{\hr\hr}w	&	{\hr\hr}ɸ	&	{\hr\hr}r	&	{\hr\hr}r	&	{\hr\hr}ʤ	&	{\hr\hr}b	&	-m	&	{\hr\hr}d	&	-n	&	{\hr\hr}ŋ	\\	
North Dobel	&	{\hr\hr}w/y/{\0}	&	{\hr\hr}ɸ	&	{\hr\hr}r	&	{\hr\hr}r	&	{\hr\hr}s	&	{\hr\hr}b	&	-m	&	{\hr\hr}d	&	-n	&	{\hr\hr}ŋ	\\	
Barakai	&	{\hr\hr}{\0}	&	{\hr\hr}w	&	{\hr\hr}r	&	{\hr\hr}r	&	{\hr\hr}ʤ	&	{\hr\hr}b	&	-m	&	{\hr\hr}d	&	-n	&	{\hr\hr}g	\\	
SW. Tarangan	&	**w > {\0}	&	{\hr\hr}ɸ	&	{\hr\hr}r	&	{\hr\hr}r	&	{\hr\hr}ʤ	&	{\hr\hr}b	&	-m	&	{\hr\hr}d	&	-n	&	{\hr\hr}ŋ	\\	
		\lspbottomrule
	\end{tabular}
		\begin{tablenotes}\tnsize
			\item[†]	{\citet{Nivens1993} states that Proto-Aru **R was “possibly a uvular trill”.}
			\item[‡]	{\citet{Nivens1993,Nivens2017} transcribes the Proto-Aru voiced plosives as *b, *d,
								and *g in accordance with his interpretation of their phonetics.
								We have retranscribed them in accordance with our interpretation.}
		\end{tablenotes}
	\end{threeparttable}
%\end{adjustbox}
\end{table}

Evidence that the Proto-Aru voiced stops were prenasalised comes
from three main sources.
Firstly, while **mb and **nd are reflected as plain voiced stops
word initially and medially, they often (though not universally) change to nasals
when they become word final due to loss of a final vowel.
Two examples are given in reflexes of **ndumbu ‘six’
and **undu ‘blood’ in \trf{06-01},
with three additional examples given in \trf{06-03}.
If Proto-Aru
**mb and **nd were prenasalised this change is simply loss of the
plosive component word finally.
Thus, for instance, **undu > **und > Kola \it{un} ‘blood’.

Secondly, Proto-Aru **ŋg > \it{ŋ} in all languages except Kola,
Barakai, and some varieties of Dobel.
One example is Proto-Aru **ŋgamáR ‘sky’ given in \trf{06-01},
with three additional examples given in \trf{06-03}.
Again, if **ŋg were prenasalised [ŋg] this is simply loss of the plosive
element in languages with **ŋg > \it{ŋ}
and loss of prenasalisation in languages with **ŋg > \it{g}.

The final piece of evidence that Proto-Aru **mb
and **nd were prenasalised comes from the fact that there are
several reconstructions in which they are derived from PMP nasal-plosive
clusters, or correspond to such a cluster in other languages.
Examples include: *ampu > **(y)ambu ‘grandparent,
grandchild’, *Rambia ‘sago palm’ > **R[ae]mbi(y)a ‘sago starch’,
*paŋdan > **ɸandan ‘pandanus’,
and Proto-South Aru **wandan ‘Banda island’.

\begin{table}\tcp{Proto-Aru prenasalised stops}{06-03}
\begin{adjustbox}{width=\textwidth}
	\begin{threeparttable}\st{2pt}
	\begin{tabular}{lllllllll}\lsptoprule
local gl.	&	wound	&	leech	&	stone	&	nose, face	&	rattan	&	rotten\\
Proto-Aru	&	**tu\tbr{mb}u	&	**i\tbr{mb}i	&	**ku\tbr{mb}u	&	**lu\tbr{ŋg}a	&	**wa\tbr{ŋg}uR	&	**ya\tbr{ŋg}al-\\ \midrule
Ujir	&	{\hr\hr}tu\tbr{b}u	&	{\hr\hr}i\tbr{b}o	&		&	{\hr\hr}lu\tbr{ŋ}i	&	{\hr\hr}wa\tbr{ŋ}i	&	\\
Kompane	&	{\hr\hr}su\tbr{b}u	&	{\hr\hr}i\tbr{b}i	&	{\hr\hr}ku\tbr{b}u	&	{\hr\hr}nu\tbr{ŋ}ay	&	{\hr\hr}gwa\tbr{ŋ}u	&	\\
East Kola	&	{\hr\hr}su\tbr{m}	&	{\hr\hr}i\tbr{m}	&		&	{\hr\hr}lu\tbr{g}i	&	{\hr\hr}wa\tbr{g}uh	&	{\hr\hr}ya\tbr{g}il\\
Lola	&		&	{\hr\hr}yi\tbr{b}i	&	{\hr\hr}ʔu\tbr{b}u	&	{\hr\hr}lu\tbr{ŋ}i	&	{\hr\hr}wa\tbr{ŋ}	&	\\
Lorang	&	{\hr\hr}u\tbr{b}u	&	{\hr\hr}i\tbr{b}i	&		&	{\hr\hr}lu\tbr{ŋ}ay	&		&	{\hr\hr}sa\tbr{ŋ}ír\\
North Dobel	&		&	{\hr\hr}yi\tbr{b}i	&		&	{\hr\hr}lu\tbr{ŋ}ay-m	&	{\hr\hr}kwa\tbr{ŋ}ur	&	{\hr\hr}sa\tbr{ŋ}al-di\\
Dobel (Warjukur)	&		&		&	{\hr\hr}ʔu\tbr{b}u	&	{\hr\hr}lu\tbr{gw}ay	&	{\hr\hr}kwa\tbr{g}ur	&	{\hr\hr}sa\tbr{g}il\\
Koba	&		&	{\hr\hr}i\tbr{b}i	&		&	{\hr\hr}nu\tbr{ŋ}ay	&	{\hr\hr}ʍa\tbr{ŋ}ar	&	{\hr\hr}sa\tbr{ŋ}il\\
Barakai	&	{\hr\hr}so\tbr{m}	&	{\hr\hr}i\tbr{b}	&		&	{\hr\hr}nu\tbr{g}e	&	{\hr\hr}go\tbr{g}or	&	{\hr\hr}ʤá\tbr{g}el\\
Karey	&		&	{\hr\hr}gí\tbr{b}ado	&		&	{\hr\hr}no\tbr{ŋ}a\su{†}	&	{\hr\hr}gwa\tbr{ŋ}ur	&	{\hr\hr}ʤá\tbr{ŋ}el\\
Batuley	&	{\hr\hr}su\tbr{m}	&		&	{\hr\hr}ku\tbr{m}	&	{\hr\hr}nu\tbr{ŋ}ei	&	{\hr\hr}gwa\tbr{ŋ}or	&	{\hr\hr}ʤá\tbr{ŋ}el\\
Mariri	&	{\hr\hr}su\tbr{m}	&		&	{\hr\hr}ko\tbr{m}	&	{\hr\hr}dʊ\tbr{ŋ}ay	&	{\hr\hr}gwa\tbr{ŋ}or	&	{\hr\hr}ʤá\tbr{ŋ}el\\
Manombai	&	{\hr\hr}su\tbr{b}u	&	{\hr\hr}ibi	&		&	{\hr\hr}lu\tbr{ŋ}ay	&		&	{\hr\hr}ʤa\tbr{ŋ}ali\\
NE. Tarangan	&		&	{\hr\hr}i\tbr{b}a	&		&	{\hr\hr}lô\tbr{ŋ}ay\su{‡}	&	{\hr\hr}gwa\tbr{ŋ}ar	&	\\
SE. Tarangan	&		&	{\hr\hr}i\tbr{b}a	&		&	{\hr\hr}lô\tbr{ŋ}ay\su{‡}	&	{\hr\hr}gwa\tbr{ŋ}ar	&	\\
NW. Tarangan	&		&	{\hr\hr}ŋi\tbr{b}ar	&		&	{\hr\hr}lu\tbr{ŋ}ay	&	{\hr\hr}ga\tbr{ŋ}ar	&	{\hr\hr}ʤa\tbr{ŋ}il\\
SW. Tarangan	&	{\hr\hr}su\tbr{m}	&	{\hr\hr}i\tbr{m}	&		&	{\hr\hr}lu\tbr{ŋ}ay	&	{\hr\hr}ga\tbr{ŋ}ar	&	{\hr\hr}ʤa\tbr{ŋ}il\\
		\lspbottomrule
	\end{tabular}
		\begin{tablenotes}\tnsize
			\item[†]	{Karey \it{noŋa} is ‘forehead’.}
			\item[‡]	{NE. and SE. Tarangan \it{lônay} ‘nose, face’ is from Gomar Meti
								and Dosi Mar villages respectively.}
		\end{tablenotes}
	\end{threeparttable}
\end{adjustbox}
\end{table}

The principal evidence for identifying the \tsc{Aru} languages
as an innovation-defined subgroup within Malayo-Polynesian is
merger of *j/*R, merger of *z/*y, and sporadic glide insertion
before certain vowel initial words.
The changes *p > **ɸ
and *b > **p can be added as supporting evidence.

	\sscp{Merger of *j/*R}{06-01-01}
\tsc{Aru} languages share a merger of PMP *j/*R.
\citet{Nivens2017} reconstructs Proto-Aru **R for the result of this merger.
Proto-Aru **R is distinct from **r (a reflex of PMP *d).
All \tsc{Aru} languages have **r = \it{r},
and all \tsc{South Aru} languages (\srf{06-02}) also have **R
> \it{r}, thus resulting in a merger of *j/*R/*d > \it{r} for that node.
However, four languages keep the reflexes of **R distinct from **r.
Ujir has **R > \it{r,
ʤ} word initially and **R > \it{y,} {\0} word medially and finally.
Kompane, \tsc{Kola}, and Lola all have **R > \it{h} word initially
and medially, and have word-final **R > \it{h,} {\0} /{\gap}{\#}.
Examples of word-initial/medial PMP *j/*R > **R are given in
\trf{06-04}, and examples in word-final position are given in \trf{06-05}.

\begin{table}[p]\centering\tcp{\tsc{Aru} *j/*R > **R /{\#\gap} and /V{\gap}V}{06-04}
\begin{adjustbox}{width=\textwidth}
	\begin{threeparttable}\st{2pt}
	\begin{tabular}{llllllll}\lsptoprule
local gl.	&	{\hr}thorn	&	{\hr}bite	&	{\hr}dry in sun	&	{\hr}nest\su{‡}	&	{\hr}name	&	{\hr}sun	&	{\hr}nose\\
PMP	&	\hr*du\tbr{R}i	&	\hr*ka\tbr{R}at	&	\hr*wa\tbr{R}i	&	\hr*\tbr{R}umaq	&	\hr*ŋa\tbr{j}an	&	\hr*qalə\tbr{j}aw	&	\hr*i\tbr{j}uŋ\\
Proto-Aru	&	**ru\tbr{R}i	&	**-ka\tbr{R}a	&	**-wa\tbr{R}ay	&	**\tbr{R}uma-	&	**ŋa\tbr{R}an	&	**la\tbr{R}aw	&	**y[ui]\tbr{R}un-\\ \midrule
Ujir	&		&	\hp{*-}\cg{(-ʤáŋ)}	&	{\hr\hr}-wo\tbr{y}	&	{\hr\hr}\tbr{ʤ}uma\su{‡}	&	{\hr\hr}ŋein	&	{\hr\hr}lau	&	{\hr\hr}ʤuin\\
Kompane	&		&	{\hr\hr}-ka\tbr{h}a	&	{\hr\hr}-gwa\tbr{h}i	&		&	{\hr\hr}ŋa\tbr{h}in	&	{\hr\hr}la\tbr{h}aw	&	{\hr\hr}ʤu\tbr{h}in\\
East Kola 	&	{\hr\hr}ri\tbr{h}i	&	{\hr\hr}-ka\tbr{h}a	&	{\hr\hr}-we\tbr{h}	&	{\hr\hr}\tbr{h}umi	&	{\hr\hr}ŋa\tbr{h}in	&	{\hr\hr}la\tbr{h}aw	&	{\hr\hr}yu\tbr{h}in\\
Lola	&		&	{\hr\hr}-ʔa\tbr{h}a	&	{\hr\hr}-wa\tbr{h}i	&		&		&	{\hr\hr}la\tbr{h}u	&	{\hr\hr}ʤu\tbr{h}in\\
Lorang	&	{\hr\hr}ru\tbr{r}	&	{\hr\hr}-ka\tbr{r}a	&	{\hr\hr}-ge\tbr{r}a	&		&		&	{\hr\hr}la\tbr{r}a	&	{\hr\hr}ʤu\tbr{r}in\\
Dobel	&		&	{\hr\hr}-ʔa\tbr{r}a	&	{\hr\hr}-kwa\tbr{r}i	&	{\hr\hr}\tbr{r}umay	&	{\hr\hr}ŋa\tbr{r}in	&	{\hr\hr}la\tbr{r}u	&	{\hr\hr}su\tbr{r}in\\
Koba	&		&	{\hr\hr}-ka\tbr{r}a	&	{\hr\hr}-ʍe\tbr{r}a	&		&		&	{\hr\hr}la\tbr{r}a	&	{\hr\hr}su\tbr{r}in\\
Barakai	&		&	\hp{*-}\cg{(r-wain)}	&	\hp{*-}\cg{(da-bálar)}	&		&	{\hr\hr}ŋa\tbr{r}an	&	{\hr\hr}lo\tbr{r}	&	{\hr\hr}ʤu\tbr{r}un\\
Karey	&	{\hr\hr}re\tbr{r}ein	&	\hp{*-}\cg{(egewain)}	&	{\hr\hr}-ɸe\tbr{r}	&		&	{\hr\hr}ŋa\tbr{r}an	&		&	{\hr\hr}ʤu\tbr{r}un\\
Batuley	&		&	{\hr\hr}-ga\tbr{r}	&	{\hr\hr}waye\tbr{r}\su{†}	&		&	{\hr\hr}ŋa\tbr{r}an	&	{\hr\hr}la\tbr{r}a	&	{\hr\hr}ʤu\tbr{r}in\\
Mariri	&		&	{\hr\hr}-ge\tbr{r}	&	{\hr\hr}-we\tbr{r}	&		&	{\hr\hr}ŋa\tbr{r}an	&	{\hr\hr}la\tbr{r}aw	&	{\hr\hr}ʤu\tbr{r}un ɸeɸ\\
Manombai	&	{\hr\hr}ru\tbr{r}i	&		&	{\hr\hr}-ga\tbr{r}ay	&		&	{\hr\hr}ŋa\tbr{r}ani	&	\hr\cg{(eran)}	&	{\hr\hr}ʤu\tbr{r}un-i\\
NE. Tar.	&	{\hr\hr}ri\tbr{r}\su{Δ}	&	{\hr\hr}-ka\tbr{r}a	&	{\hr\hr}-gwa\tbr{r}ɛ	&		&	{\hr\hr}ŋa\tbr{r}in	&		&	\\
SE. Tarang.	&	{\hr\hr}ra\tbr{r}ír-da\su{Δ}	&	{\hr\hr}-ga\tbr{r}	&	{\hr\hr}-gwa\tbr{r}ay	&		&	{\hr\hr}ŋa\tbr{r}an	&		&	{\hr\hr}ʤɔ\tbr{r}in\\
NW. Tar.	&	{\hr\hr}ri\tbr{r}ír	&	{\hr\hr}-ka\tbr{r}a	&	{\hr\hr}-ga\tbr{r}ay	&		&	{\hr\hr}ŋa\tbr{r}an	&	\hr\cg{(ŋɛl)}	&	{\hr\hr}ʤu\tbr{r}in\\
SW. Tar.	&	{\hr\hr}urú\tbr{r}-di	&	{\hr\hr}-ka\tbr{r}a	&	{\hr\hr}-ga\tbr{r}ay	&	{\hr\hr}\tbr{r}umay	&	{\hr\hr}ŋa\tbr{r}an	&	\hr{(ŋɛl)}	&	{\hr\hr}ʤu\tbr{r}in\\
		\lspbottomrule
	\end{tabular}
		\begin{tablenotes}\tnsize
			\item[†]	{Batuley \it{wayer} ‘dry in sun’ has metathesis
								of the medial and final consonants from **warey.}
			\item[‡]	{PMP *Rumaq is reconstructed with the meaning ‘house’.
								This meaning is retained by Ujir \it{ʤuma}.
								Other \tsc{Aru} reflexes have shifted to the semantics ‘nest, den’.}
			\item[Δ]	{Southeast Tarangan \it{rarír-da} ‘thorn’ is from Dosi Mar village
								and Northeast Tarangan \it{rir} ‘thorn’ is from Gomar Meti village,
								the latter supplied by Rick Nivens (pers. comm. April 2024).}
		\end{tablenotes}
	\end{threeparttable}
\end{adjustbox}
\end{table}

\begin{table}[p]
	\centering\tcp{\tsc{Aru} *j/*R > **R /{\gap\#}}{06-05}
	\st{5.5pt}\begin{tabular}{llllll}\lsptoprule
local gl.	&	{\hr}egg	&	{\hr}coconut	&	{\hr}hear	&	{\hr}fly (n.)	&	{\hr}caterpillar\\
PMP	&	\hr*qatəlu\tbr{R}	&	\hr*niu\tbr{R}	&	\hr*dəŋə\tbr{R}	&	\hr*lalə\tbr{j}	&	\hr*qulə\tbr{j}\\
Proto-Aru	&	**tulu\tbr{R}-	&	**nu\tbr{R}	&	**-reŋa\tbr{R}	&	**l[ae]l[iu]\tbr{R}	&	**[wy]ila\tbr{R}\\ \midrule
Ujir	&	{\hr\hr}tuli	&	{\hr\hr}nu\tbr{y}	&	{\hr\hr}-reŋa\tbr{y}	&		&	{\hr\hr}wil{\tl}wila\tbr{y}\\
		&		&		&		&		&	{\hr\hr}tun{\tl}tun\\
Kompane	&	{\hr\hr}tuli	&	{\hr\hr}nu\tbr{h}	&	{\hr\hr}-reŋa\tbr{h}	&		&	\\
East Kola	&	{\hr\hr}tuli\tbr{h}	&	{\hr\hr}nu\tbr{h}	&	{\hr\hr}-rena	&	{\hr\hr}lelu\tbr{h}	&	{\hr\hr}yila\tbr{h}\\
Dobel	&	{\hr\hr}tuli\tbr{r}	&	{\hr\hr}no\tbr{r}	&	{\hr\hr}-reŋin	&	{\hr\hr}lali\tbr{r}	&	{\hr\hr}s{\tl}sila\tbr{r}\\
Manombai	&	{\hr\hr}tulu\tbr{r}-i	&	{\hr\hr}no\tbr{r}	&	{\hr\hr}-reŋa\tbr{r}	&		&	\\
Batuley	&	{\hr\hr}tuli\tbr{r}	&	{\hr\hr}no\tbr{r}	&	{\hr\hr}-reŋa\tbr{r}	&	{\hr\hr}lale\tbr{r}	&	\\
SW. Tarangan	&	{\hr\hr}tuli\tbr{r}	&	{\hr\hr}nô\tbr{r}	&		&		&	{\hr\hr}ʤila\tbr{r}\\
		\lspbottomrule
	\end{tabular}
\end{table}

	\sscp{Merger of *z/*y}{06-01-02}
Another innovation shared by all Aru languages is merger of PMP
*z/*y > **y.\footnote{
Merger of *z/*y > **y also occurs in \tsc{Sula-Buru}.
This is probably a parallel development.}
This change also affects
PMP *i before vowels (/{\gap}V),
including when it occurs before a vowel due to loss of *h or *q.
\citet{Nivens1993,Nivens2017} reconstructs Proto-Aru **y as the
result of this merger.
Word initially, most languages then have fortition of **y > **ʤ,
with Dobel and Koba undergoing subsequent **ʤ > \it{s}.
East Kola retains initial **y = \it{y} in most cases,
but some lexemes have \it{y} {\tl} \it{ʤ} in variation.
Examples of PMP *z/*y/*i > **y in word-initial position are given
in \trf{06-06}.

\begin{table}\tcp{\tsc{Aru} *z/*y/*i > **y /{\#\gap}}{06-06}
%\begin{adjustbox}{width=\textwidth}
	\begin{threeparttable}\st{5pt}
	\begin{tabular}{lllllll}\lsptoprule
local gl.	&	far	&	path	&	sew	&	blow	&	count, speak	&	shark\\
PMP	&	\hr*\tbr{z}auq	&	\hr*\tbr{z}alan	&	\hr*\tbr{z}aqit	&	\hr*hə\tbr{y}up	&	\hr*\tbr{i}hap	&	\hr*q\tbr{i}hu\\
Proto-Aru	&	**\tbr{y}aw-	&	**\tbr{y}ala	&	**-\tbr{y}ay	&	**-\tbr{y}uɸ	&	**-\tbr{y}aɸ	&	**\tbr{y}u\\ \midrule
Ujir	&	{\hr\hr}\tbr{ʤ}aw	&	{\hr\hr}\tbr{ʤ}ala	&	{\hr\hr}-\tbr{ʤ}e\tbr{y}	&	{\hr\hr}-\tbr{ʤ}u	&		&	{\hr\hr}\tbr{ʤ}u\\
Kompane	&	{\hr\hr}\tbr{ʤ}aw	&	{\hr\hr}\tbr{ʤ}ala	&		&		&		&	\\
West Kola	&	{\hr\hr}\tbr{ʤ}aw	&	{\hr\hr}\tbr{ʤ}ala	&		&		&		&	\\
East Kola	&	{\hr\hr}\tbr{y}aw	&	{\hr\hr}\tbr{y}ala	&	{\hr\hr}-\tbr{y}e\tbr{y}	&	{\hr\hr}-\tbr{y}uf	&	{\hr\hr}-\tbr{y}af	&	{\hr\hr}\tbr{y}u\\
Lola	&	{\hr\hr}\tbr{ʤ}aw	&	{\hr\hr}\tbr{ʤ}ala	&		&		&		&	\\
Lorang	&	{\hr\hr}\tbr{ʤ}aw	&	{\hr\hr}\tbr{ʤ}ala	&		&		&		&	\\
Dobel	&	{\hr\hr}\tbr{s}aw	&	{\hr\hr}\tbr{s}ala	&	\hr\cg{(-sela)}	&	{\hr\hr}-\tbr{s}o	&	{\hr\hr}-\tbr{s}aw	&	{\hr\hr}\tbr{s}u\\
Koba	&	{\hr\hr}\tbr{s}aw	&	{\hr\hr}\tbr{s}ala	&		&		&		&	\\
Barakai	&	{\hr\hr}\tbr{ʤ}aw	&		&		&		&		&	\\
Karey	&	{\hr\hr}\tbr{ʤ}aw	&	{\hr\hr}\tbr{ʤ}al	&		&		&		&	\\
Batuley	&	{\hr\hr}\tbr{ʤ}a	&		&		&		&		&	{\hr\hr}\tbr{ʤ}u\\
Mariri	&	{\hr\hr}\tbr{ʤ}aw	&	{\hr\hr}\tbr{ʤ}el	&		&		&		&	\\
Manombai	&	{\hr\hr}\tbr{ʤ}aw	&	{\hr\hr}\tbr{ʤ}ala	&		&	{\hr\hr}-\tbr{ʤ}u	&		&	\\
NE. Tarangan	&	{\hr\hr}\tbr{ʤ}ɔ	&	{\hr\hr}\tbr{ʤ}ala	&		&		&		&	\\
SE. Tarangan	&	{\hr\hr}\tbr{ʤ}aw	&	{\hr\hr}\tbr{ʤ}al	&		&		&		&	\\
NW. Tarangan	&	{\hr\hr}\tbr{ʤ}ɔ	&	{\hr\hr}\tbr{ʤ}ala	&		&	{\hr\hr}-\tbr{ʤ}ô\su{†}	&	{\hr\hr}-\tbr{ʤ}ɔ\su{†}	&	{\hr\hr}\tbr{ʤ}u\\
SW. Tarangan	&	{\hr\hr}\tbr{ʤ}ôw	&	{\hr\hr}\tbr{ʤ}ala	&	{\hr\hr}-\tbr{ʤ}ɛ\tbr{y}	&	{\hr\hr}-\tbr{ʤ}ôw	&	{\hr\hr}-\tbr{ʤ}ôw	&	{\hr\hr}\tbr{ʤ}u\\
		\lspbottomrule
	\end{tabular}
		\begin{tablenotes}\tnsize
			\item[†] {Northwest Tarangan \it{-ʤô} ‘blow’ and \it{-}ʤɔ ‘count, speak’ are from Doka Timur village.}
		\end{tablenotes}
	\end{threeparttable}
%\end{adjustbox}
\end{table}

The reflexes of *z/*y > **y are more complex when word medial,
with many languages having an unconditioned split.
The most common reflexes of medial **y are \it{y,} {\0},
or loss with fronting of an adjacent vowel \tsc{(+fr)}.
Some languages have glide fortition with **y > **ʤ,
and Dobel and Koba have subsequent **ʤ > \it{s}.
The reflexes of *z/*y/*i > **y in word-medial position attested
in each language are summarised in \trf{06-07}, and examples are
given in \trf{06-08}.

\begin{table}
	\centering\tcp{Reflexes of PMP *z/*y/*i word medially in \tsc{Aru}}{06-07}
	\st{6pt}\begin{tabular}{llllll||l|lllll}\lsptoprule
Ujir	&	y	&		&		&	\tsc{+fr}	&	{\0}	&	Karey	&		&		&		&		&	{\0}	\\
Kompane	&		&		&		&	\tsc{+fr}	&	{\0}	&	Batuley	&	y	&		&		&	\tsc{+fr}	&	{\0}	\\
West Kola	&		&		&		&		&	{\0}	&	Mariri	&	y	&	ʤ	&		&	\tsc{+fr}	&		\\
East Kola	&		&		&		&	\tsc{+fr}	&	{\0}	&	Manombai	&	y	&		&		&	\tsc{+fr}	&		\\
Lola	&	y	&		&		&	\tsc{+fr}	&	{\0}	&	NE. Tarangan	&	y	&		&		&		&	{\0}	\\
Lorang	&		&	ʤ	&		&		&	{\0}	&	SE. Tarangan	&		&	ʤ	&		&		&	{\0}	\\
Dobel	&		&		&	s	&		&	{\0}	&	NW. Tarangan	&	y	&		&		&	\tsc{+fr}	&	{\0}	\\
Koba	&		&		&	s	&		&	{\0}	&	SW. Tarangan	&	y	&		&		&		&	{\0}	\\
Barakai	&		&		&		&		&	{\0}	&		&		&		&		&		&		\\
		\lspbottomrule
	\end{tabular}
\end{table}

\begin{table}\tcp{\tsc{Aru} *z/*y/*i > **y /V{\gap}V}{06-08}
\begin{adjustbox}{width=\textwidth}
	\begin{threeparttable}\st{2pt}
	\begin{tabular}{llllllll}\lsptoprule
local gl.	&	{\hr}rain	&	{\hr}ladder	&	\hp{-}{\hr}pay	&	{\hr}fresh water	&	{\hr}mango	&	{\hr}interior	&	{\hr}dugong\\
PMP	&	\hr*qu\tbr{z}an	&	\hr*haRə\tbr{z}an	&	\hp{-}\hr*ba\tbr{y}aD	&	\hr*wahiyəR	&	\hr*wa\tbr{i}	&	\hr*da\tbr{y}a	&	\hr*du\tbr{y}uŋ\\
Proto-Aru	&	**\cg{(w)}u\tbr{y}an	&	**Re\tbr{y}an	&	**-pa\tbr{y}aR	&	**wa\tbr{y}aR	&	**wa\tbr{y}[a]	&	**ra\tbr{y}a\su{†}	&	**r\cg{(iu)}\tbr{y}u\\ \midrule
Ujir	&	{\hr\hr}uan	&	{\hr\hr}rián	&	{\hr\hr}-ɸa\tbr{y}	&	{\hr\hr}wa\tbr{y}	&	{\hr\hr}wa	&	{\hr\hr}raa	&	{\hr\hr}riu, ri\tbr{y}u\\
Kompane	&	{\hr\hr}gwian	&		&		&	{\hr\hr}gwah	&		&	{\hr\hr}raa	&	\\
West Kola	&	{\hr\hr}uan	&		&		&	{\hr\hr}wah	&		&		&	\\
East Kola	&	{\hr\hr}uin	&	{\hr\hr}heen	&	{\hr\hr}-ɸeeh	&	{\hr\hr}weeh	&	{\hr\hr}wee	&	{\hr\hr}ree	&	{\hr\hr}riu, ri\tbr{y}u\\
Lola	&	{\hr\hr}wian	&		&		&	{\hr\hr}waʔ	&		&	{\hr\hr}ra\tbr{y}a	&	\\
Lorang	&	{\hr\hr}gú\tbr{ʤ}en	&		&		&	{\hr\hr}gwar	&		&	\hr\cg{(wúney)}	&	\\
Dobel	&	{\hr\hr}ʔu\tbr{s}an	&		&	{\hr\hr}-ɸa\tbr{s}ir	&	{\hr\hr}kwar	&	{\hr\hr}kwa\tbr{s}i	&	{\hr\hr}ra\tbr{s}a	&	{\hr\hr}riw\\
Koba	&	{\hr\hr}ʔu\tbr{s}an	&		&		&	{\hr\hr}ʍar	&		&	\hr\cg{(wúney)}	&	\\
Barakai	&	{\hr\hr}guon	&		&		&	{\hr\hr}gwaor	&		&		&	\\
Karey	&	{\hr\hr}guon	&		&		&	{\hr\hr}gaor	&		&		&	\\
Batuley	&	{\hr\hr}guon	&	{\hr\hr}rion	&		&	{\hr\hr}gwa\tbr{y}or	&		&	{\hr\hr}ra\tbr{y}	&	{\hr\hr}ri\\
Mariri	&	{\hr\hr}gu\tbr{ʤ}on	&		&		&	{\hr\hr}gwa\tbr{y}or	&		&	{\hr\hr}ree	&	\\
Manombai	&	{\hr\hr}gian	&	{\hr\hr}rean	&		&	{\hr\hr}g\tbr{a}yar	&		&	{\hr\hr}rea	&	\\
NE. Tarangan	&	{\hr\hr}gu\tbr{y}an	&		&		&	{\hr\hr}gwar	&		&	{\hr\hr}ra	&	\\
SE. Tarangan	&	{\hr\hr}gu\tbr{ʤ}an	&		&		&	{\hr\hr}gwar	&		&	{\hr\hr}ra	&	\\
NW. Tarangan	&	{\hr\hr}gian	&	{\hr\hr}rêan\su{‡}	&		&	{\hr\hr}gar	&	{\hr\hr}ga\tbr{y}a\su{‡}	&	{\hr\hr}ra	&	{\hr\hr}ru\\
SW. Tarangan	&	{\hr\hr}gô\tbr{y}an	&	{\hr\hr}re\tbr{y}an	&	{\hr\hr}-ɸa\tbr{y}ar	&	{\hr\hr}gar	&	{\hr\hr}ga\tbr{y}a	&	{\hr\hr}ra	&	{\hr\hr}ru\\
		\lspbottomrule
	\end{tabular}
		\begin{tablenotes}\tnsize
			\item[†]	{PMP *daya is reconstructed with the meaning ‘upriver, toward the interior’.
								Reflexes in the Aru languages mean ‘wooded area; land/interior’.
								East Kola \it{ree} has these meanings in addition to ‘weather’.}
			\item[‡]	{Northwest Tarangan \it{rêan} ‘ladder’
								and \it{gaya} ‘mango’ are from Rebi village.}
		\end{tablenotes}
	\end{threeparttable}
\end{adjustbox}
\end{table}

	\sscp{Initial glide insertion}{06-01-03}
Another change which can be assigned to the level of Proto-Aru
is sporadic insertion of initial glides before vowel initial
words (including words which became vowel initial due to loss
of PMP *q and *h), shown in \trf{06-09}.
In most cases the glide inserted is palatal *y,
but *w is often inserted when the word contains *u.


\begin{table}\tcp{\tsc{Aru} word-initial glide insertion}{06-09}
\begin{adjustbox}{width=\textwidth}
	\begin{threeparttable}\st{2pt}
	\begin{tabular}{lllllll}\lsptoprule
local gl.	&	{\hr}nose	&	{\hr}tail	&	{\hr}deep sea\su{‡}	&	{\hr}fish	&	{\hr}grass	&	{\hr}rain\\
PMP	&	\hr*ijuŋ	&	\hr*ikuR	&	\hr*qaluR	&	\hr*hikan	&	\hr*udu	&	\hr*quzan\\
Proto-Aru	&	**\tbr{y}uR[ui]n-	&	**\tbr{y}uk[ui]R-	&	**\tbr{w}aluR	&	**(\tbr{y})ika	&	**(\tbr{w})uru	&	**(\tbr{w})uyan\\ \midrule
Ujir	&	{\hr\hr}\tbr{ʤ}uin	&		&	{\hr\hr}\tbr{w}ali	&	{\hr\hr}ika	&	{\hr\hr}ur	&	{\hr\hr}uan\\
Kompane	&	{\hr\hr}\tbr{ʤ}uhin	&	{\hr\hr}\tbr{ʤ}ukih	&		&	{\hr\hr}\tbr{ʤ}ika	&		&	{\hr\hr}\tbr{gw}ian\\
West Kola	&	{\hr\hr}\tbr{y}uhin	&	{\hr\hr}\tbr{ʤ}ukih	&		&	{\hr\hr}ika	&		&	{\hr\hr}uan\\
East Kola	&	{\hr\hr}\tbr{y}uhin	&	{\hr\hr}\tbr{y}ukih	&	{\hr\hr}\tbr{w}aluh	&	{\hr\hr}ika	&	{\hr\hr}ur	&	{\hr\hr}uin\\
Lola	&	{\hr\hr}\tbr{ʤ}uhin	&	{\hr\hr}\tbr{ʤ}uʔi	&		&	{\hr\hr}\tbr{ʤ}iʔa	&		&	{\hr\hr}\tbr{w}ian\\ \midrule
Proto-S. Aru	&	**\tbr{y}ur[ui]n-	&	**\tbr{y}uk[ui]R-	&	**\tbr{w}alur	&	**\tbr{y}ika	&	**\tbr{w}uru	&	**\tbr{w}uyan\\ \midrule
Lorang	&	{\hr\hr}\tbr{ʤ}urin	&	{\hr\hr}\tbr{ʤ}ukir	&		&		&		&	{\hr\hr}\tbr{g}úʤen\\
Dobel	&	{\hr\hr}\tbr{s}urin	&	{\hr\hr}\tbr{s}uʔir	&	{\hr\hr}\tbr{kw}alur	&	{\hr\hr}\tbr{s}iʔa	&	{\hr\hr}\tbr{k}uru	&	{\hr\hr}\tbr{k}usan\\
Koba	&	{\hr\hr}\tbr{s}urin	&	{\hr\hr}\tbr{s}ukir	&		&	{\hr\hr}\tbr{s}ika	&		&	{\hr\hr}usan\\
Barakai	&	{\hr\hr}\tbr{ʤ}urun	&	{\hr\hr}\tbr{ʤ}or	&		&	{\hr\hr}\tbr{ʤ}ɪg	&		&	{\hr\hr}\tbr{g}uon\\
Karey	&	{\hr\hr}\tbr{ʤ}urun	&	{\hr\hr}\tbr{ʤ}úor	&		&	{\hr\hr}\tbr{ʤ}ego	&		&	{\hr\hr}\tbr{g}uon\\
Batuley	&	{\hr\hr}\tbr{ʤ}urin	&	{\hr\hr}\tbr{ʤ}ugor	&	{\hr\hr}\tbr{gw}alor	&	{\hr\hr}\tbr{ʤ}ig	&	{\hr\hr}\tbr{g}ur	&	{\hr\hr}\tbr{g}uon\\
Mariri	&	{\hr\hr}\tbr{ʤ}urun ɸeɸ	&	{\hr\hr}\tbr{ʤ}úgor	&		&	{\hr\hr}\tbr{ʤ}ɪg	&		&	{\hr\hr}\tbr{g}uʤon\\
Manombai	&	{\hr\hr}\tbr{ʤ}urun-i	&		&		&		&		&	{\hr\hr}\tbr{g}ian\\
NE. Tarangan	&		&	{\hr\hr}\tbr{ʤ}ukir	&	{\hr\hr}\tbr{gw}alɔ	&	{\hr\hr}\tbr{ʤ}ɛka	&		&	{\hr\hr}\tbr{g}uyan\\
SE. Tarangan	&	{\hr\hr}\tbr{ʤ}ɔrin	&	{\hr\hr}\tbr{ʤ}ugir	&	{\hr\hr}\tbr{gw}alɔr\su{†}	&	{\hr\hr}\tbr{ʤ}egɔ	&		&	{\hr\hr}\tbr{g}uʤan\\
NW. Tarangan	&	{\hr\hr}\tbr{ʤ}urin	&		&		&		&	{\hr\hr}\tbr{g}ur	&	{\hr\hr}\tbr{g}ian\\
SW. Tarangan	&	{\hr\hr}\tbr{ʤ}urin	&		&	{\hr\hr}\tbr{g}alar	&		&		&	{\hr\hr}\tbr{g}ôyan\\
		\lspbottomrule
	\end{tabular}
		\begin{tablenotes}\tnsize
			\item[†]	{Southeast Tarangan \it{gwalɔr} ‘deep sea’ is from Beltubur village.}
			\item[‡]	{East Kola \it{waluh} is ‘deep water area of sea’,
								Dobel \it{kwalur} is ‘a specific (named) area of deep sea’,
								East Tarangan \it{gwalɔ} is ‘saltwater,
								sea’, Batuley \it{gwalor} is ‘deep blue sea’,
								Southwest Tarangan \it{galar} is ‘deep water passage/channel’
								and is from Feruni village.}
		\end{tablenotes}
	\end{threeparttable}
\end{adjustbox}
\end{table}

Similar insertions are fairly widespread in the region (see \srf{13-02-02}),
meaning that it constitutes only weak subgrouping evidence.
Furthermore, in \tsc{Aru} glide insertion is sporadic
and does not affect all vowel initial words.
Nonetheless, there are several words for which all reflexes attest
an initial glide, thus allowing reconstruction of the glide to Proto-Aru
for those etyma.
It thus can be taken as supporting evidence for the validity of Proto-Aru.
Three examples are *ijuŋ > **y[ui]Run- ‘nose’,
*ikuR > **y[ui]kuR- ‘tail’,
and *qaluR > **aluR > **waluR ‘deep sea’,
shown in \trf{06-09}.

After glide insertion at the level of Proto-Aru there appears
to have been additional rounds of glide insertion.
One had occurred by the time of Proto-South Aru (see \srf{06-02}),
leading \citet{Nivens2017} to reconstruct an unambiguous glide
for Proto-South Aru in some words,
as seen in *hikan > **(y)ika ‘fish’ > Proto-South Aru **yika
‘fish’, *udu > **(w)uru > **wuru ‘grass,
weeds’, and *quzan > **(w)uyan > **wuyan ‘rain’,
shown in \trf{06-09}.
An additional round of glide insertion also
affected Lola, Dobel, Koba, and Lorang.
This last round of glide insertion occurred after the fortition
of glides previously inserted in Dobel,
Koba, and Lorang, thus reintroducing word-initial glides into these languages.

The data regarding glide insertion in the \tsc{Aru} languages
is probably best explained by positing that it was a change in
progress which had begun at the level of Proto-Aru
and continued to varying degrees in some,
but not all, groups after the break-up of Proto-Aru.

	\largerpage
\sscp{Labial plosives}{06-01-04}
Additional evidence for Proto-Aru comes from the sound changes
affecting the labial plosives.
We begin with the changes affecting PMP *p.
The changes affecting PMP *p are simplest in Ujir,
Kompane, \tsc{Kola}, and Lola.
Ujir has PMP *p > {\0} initially
and PMP *p > \it{w,} {\0} medially.
\tsc{Kola} has PMP *p > \it{f} in all positions.
Kompane has PMP *p > \it{h} initially,
and PMP *p > \it{h,} {\0} word medially.
Lola has initial PMP *p > \it{h},
but medial PMP *p > {\0},
\it{w, h,} or rounding of an adjacent vowel \tsc{(+round)}.
The changes affecting PMP *p in the \tsc{South Aru} languages
are more complex, with several splits attested.
The different reflexes are summarised in \trf{06-10}.
They include: PMP *p > \it{w},
fortition (of intermediate **w),
loss, and rounding of an adjacent vowel.
\citet{Nivens2017} posits Proto-South Aru **β to account for these reflexes.
Examples of PMP *p in Aru are given in \trf{06-11}.

\begin{table}
	\centering\tcp{Reflexes of PMP *p in \tsc{South Aru}}{06-10}
	\st{6pt}\begin{tabular}{llllll}\lsptoprule
Lorang	&		&	w	&	gw	&		&	{\0}\\
Dobel	&		&	w	&	kw	&		&	{\0}\\
Koba	&		&	o	&	ʍ	&		&	{\0}\\
Barakai	&		&		&		&		&	{\0}\\
Karey	&		&	β	&	g	&	\tsc{+round}	&	{\0}\\
Batuley	&		&	w	&		&		&	{\0}\\
Mariri	&		&	w	&		&		&	{\0}\\
Manombai	&		&		&		&		&	{\0}\\
\tsc{Tarangan}	&		&		&		&	\tsc{+round}	&	{\0}\\
		\lspbottomrule
	\end{tabular}
\end{table}

\begin{table}\tcp{\tsc{Aru} *p > **ɸ}{06-11}
\begin{adjustbox}{width=\textwidth}
	\begin{threeparttable}\st{2pt}
	\begin{tabular}{llllllll}\lsptoprule
local gl.	&	{\hr}fire	&	{\hr}bat	&	{\hr}thin	&	{\hr}hungry	&	{\hr}tooth	&	{\hr}shoot	&	{\hr}ten\\
PMP	&	\hr*ha\tbr{p}uy	&	\hr*\tbr{p}aniki	&	\hr*ma-ni\tbr{p}is	&	\hr*la\tbr{p}aR	&	\hr*[nŋ]i\tbr{p}ən	&	\hr*\tbr{p}anaq	&	\\
Proto-Aru	&	**a\tbr{ɸ}uy	&	**\tbr{ɸ}aniki	&	**mani\tbr{ɸ}i-	&	**kan-la\tbr{ɸ}aR	&	**[nŋ]e\tbr{ɸ}in	&	**-\tbr{ɸ}ana	&	**\tbr{ɸ}uR\su{Δ}\\ \midrule
Ujir	&	{\hr\hr}ʤe	&		&	{\hr\hr}maníʔ	&	{\hr\hr}kana\tbr{w}ay	&	{\hr\hr}nein	&	{\hr\hr}-an	&	{\hr\hr}uy sia\\
Kompane	&	{\hr\hr}a\tbr{h}i	&	{\hr\hr}\tbr{h}anisi	&	{\hr\hr}maniay	&	{\hr\hr}kanaa	&	{\hr\hr}ne\tbr{h}in	&		&	{\hr\hr}\tbr{h}uɸaɸi\\
East Kola	&	{\hr\hr}e\tbr{f}	&	{\hr\hr}\tbr{f}anís	&	{\hr\hr}mani\tbr{f}i	&	{\hr\hr}kanaah-ni	&	{\hr\hr}no\tbr{f}in	&	{\hr\hr}-\tbr{f}an	&	{\hr\hr}\tbr{f}uh\\
Lola	&	{\hr\hr}ʤa\tbr{h}i	&	{\hr\hr}\tbr{h}anisi	&	{\hr\hr}manuay	&	{\hr\hr}ʔaná\tbr{w}i	&	{\hr\hr}ŋein	&		&	{\hr\hr}\tbr{h}aɸi ye\\  \midrule
Lorang	&	{\hr\hr}ʤay	&	{\hr\hr}\tbr{w}eníu	&	{\hr\hr}mani\tbr{gw}ay	&		&	{\hr\hr}yɪŋɪn	&		&	{\hr\hr}uru\\
Dobel	&	{\hr\hr}say	&	{\hr\hr}\tbr{w}oyníw	&	{\hr\hr}mani\tbr{kw}ay	&	{\hr\hr}na-ka‹l›la\tbr{w}ar-ni	&	{\hr\hr}ŋeyan-	&	{\hr\hr}-\tbr{w}ana	&	{\hr\hr}\tbr{w}ur\\
Koba	&	{\hr\hr}say	&	{\hr\hr}\tbr{o}éni	&	{\hr\hr}maní\tbr{ʍ}ey	&		&		&		&	{\hr\hr}ur\\
Barakai	&	{\hr\hr}aa	&	{\hr\hr}aaní	&	{\hr\hr}mania	&	{\hr\hr}kanáir	&	{\hr\hr}neen	&		&	{\hr\hr}urwéo\\
Karey	&	{\hr\hr}aa	&	{\hr\hr}anido	&	{\hr\hr}manua	&	{\hr\hr}kaná\tbr{β}ir	&	{\hr\hr}ne\tbr{g}en	&		&	{\hr\hr}orwéw\\
Batuley	&	{\hr\hr}ʤekáy	&	{\hr\hr}anís	&	{\hr\hr}mani\tbr{w}e	&	{\hr\hr}kana\tbr{w}er	&	{\hr\hr}neyen	&		&	{\hr\hr}ur\\
Mariri	&	{\hr\hr}ay	&	{\hr\hr}adíʧ	&	{\hr\hr}madiway	&		&	{\hr\hr}deen	&		&	{\hr\hr}orɸéɸ\\
Manombai	&	{\hr\hr}a\tbr{w}	&	{\hr\hr}nisi	&	{\hr\hr}maniay	&	{\hr\hr}kanári	&	{\hr\hr}nini	&	{\hr\hr}-en\su{‡}	&	{\hr\hr}uraɸa ia\\
NE. Tarang.	&	{\hr\hr}ɛ	&	{\hr\hr}ɔnít	&	{\hr\hr}manuɛ	&		&	{\hr\hr}nin	&		&	{\hr\hr}urɸeɸa\\
SE. Tarang.	&	{\hr\hr}ɛy	&	{\hr\hr}ɔnít	&	{\hr\hr}mɛnugway	&	{\hr\hr}kanagwir	&	{\hr\hr}nɛin	&		&	{\hr\hr}ɔrɸɛɸɔ\\
NW. Tar.	&	{\hr\hr}ɔ	&	{\hr\hr}nik	&	{\hr\hr}manuay	&	{\hr\hr}kanɔir	&	{\hr\hr}nin	&	{\hr\hr}-ɔn	&	{\hr\hr}ur\\
SW. Tar.	&	{\hr\hr}ɔ\tbr{w}	&	{\hr\hr}aník\su{†}	&	{\hr\hr}mɛnuay	&	{\hr\hr}kanɔir	&	{\hr\hr}nin	&	{\hr\hr}-ɔn	&	{\hr\hr}ur\\
		\lspbottomrule
	\end{tabular}
		\begin{tablenotes}\tnsize
			\item[†]	{Southwest Tarangan \it{aník} ‘fruit bat’ is from Feruni and Kabalukin villages.}
			\item[‡]	{Manombai \it{-en} ‘shoot’ is from Gardakau village.}
			\item[Δ]	{Several reflexes of **ɸuR ‘ten’ here are from a compound **ɸuR-papay.}
		\end{tablenotes}
	\end{threeparttable}
\end{adjustbox}
\end{table}

PMP *b > \it{ɸ} in most Aru languages.
\tsc{Tarangan} usually has *b > \it{p} when it becomes word final
due to loss of an earlier vowel,
and *b > \it{ɸ {\tl} ʍ {\tl} h} in other word positions with
the result varying by the variety of \tsc{Tarangan} and/or by
quality of the following vowels.
Karey has *b > ɸ /{\gap}o
and *b > \it{w} before other vowels.
Barakai has *b > \it{w}.

While Kola has *b > \it{ɸ} in all word positions the phoneme /ɸ/ has allophony.
\citet{Nivens2017} describes Kola /ɸ/ as “(always?) rounded,
and sometimes half-voiced”, while \citet[34]{TakataTakata1992}
describe free variation between a voiceless
and voiced bilabial fricative,
/ɸ/ → [ɸ] {\tl} [β].
We transcribe this segment ‹ɸʷ›.
Kola bilabial /ɸ/ contrasts with labio-dental /f/.

\citet{Nivens2017} reconstructs the reflex of PMP *b in Proto-Aru
as **p, noting that an alternate reconstruction would be **β.
Examples of PMP *b > **p are given in \trf{06-12}.

\begin{table}\tcp{\tsc{Aru} *b > **p}{06-12}
\begin{adjustbox}{width=\textwidth}
	\begin{threeparttable}\st{2pt}
	\begin{tabular}{llllllll}\lsptoprule
local gl.	&	{\hr}body hair	&	{\hr}moon	&	{\hr}fruit	&	{\hr}stone	&	{\hr}mouth	&	{\hr}top	&	{\hr}pig\\
PMP	&	\hr*\tbr{b}ulu	&	\hr*\tbr{b}ulan	&	\hr*\tbr{b}uaq	&	\hr*\tbr{b}atu	&	\hr*\tbr{b}aq\tbr{b}aq	&	\hr*\tbr{b}u\tbr{b}uŋ	&	\hr*\tbr{b}a\tbr{b}uy\\
Proto-Aru	&	**\tbr{p}ulu-	&	**\tbr{p}ulan	&	**\tbr{p}ua-\su{‡}	&	**\tbr{p}atu	&	**\tbr{p}a\tbr{p}a-	&	**\tbr{p}u\tbr{p}un	&	**\tbr{p}a\tbr{p}uy\\ \midrule
Ujir	&	{\hr\hr}\tbr{ɸ}ul	&	{\hr\hr}\tbr{ɸ}ulan	&	{\hr\hr}\tbr{ɸ}uay	&	{\hr\hr}\tbr{ɸ}at	&	{\hr\hr}\tbr{ɸ}e\tbr{ɸ}a	&		&	\\
Kompane	&		&	{\hr\hr}\tbr{ɸ}ulan	&	{\hr\hr}\tbr{ɸ}ui	&		&	{\hr\hr}\tbr{ɸ}a\tbr{ɸ}i	&		&	\\
East Kola	&		&	{\hr\hr}\tbr{ɸʷ}ulan	&	{\hr\hr}\tbr{ɸʷ}uui	&	{\hr\hr}\tbr{ɸʷ}at	&	{\hr\hr}\tbr{ɸʷ}a\tbr{ɸʷ}a-	&	{\hr\hr}\tbr{ɸʷ}u\tbr{ɸʷ}un\su{¶}	&	{\hr\hr}\tbr{ɸʷ}e\tbr{ɸʷ}\\
Lola	&	{\hr\hr}\tbr{ɸ}uli	&	{\hr\hr}\tbr{ɸ}ulan	&	{\hr\hr}\tbr{ɸ}ua-di	&		&	{\hr\hr}\tbr{ɸ}a\tbr{ɸ}i	&	{\hr\hr}\tbr{ɸ}u\tbr{ɸ}in	&	\\
Lorang	&	{\hr\hr}\tbr{ɸ}ul	&	{\hr\hr}\tbr{ɸ}ulan	&	{\hr\hr}\tbr{ɸ}úʤa	&	{\hr\hr}\tbr{ɸ}atu	&	{\hr\hr}\tbr{ɸ}a\tbr{ɸ}a	&		&	\\
Dobel	&	{\hr\hr}\tbr{ɸ}ulu-\su{ǁ}	&	{\hr\hr}\tbr{ɸ}ulan	&	{\hr\hr}\tbr{ɸ}usi	&	{\hr\hr}\tbr{ɸ}atu	&	{\hr\hr}\tbr{ɸ}a\tbr{ɸ}a-m	&	{\hr\hr}\tbr{ɸ}u\tbr{ɸ}un	&	{\hr\hr}\tbr{ɸ}a\tbr{ɸ}i\\
Koba	&	{\hr\hr}\tbr{ɸ}ul	&	{\hr\hr}\tbr{ɸ}ulan	&	{\hr\hr}\tbr{ɸ}úsa	&	{\hr\hr}\tbr{ɸ}atu	&	{\hr\hr}\tbr{ɸ}e\tbr{ɸ}a	&	{\hr\hr}\tbr{ɸ}i\tbr{ɸ}íŋ	&	\\
Barakai	&	{\hr\hr}\tbr{w}olʊ-n	&	{\hr\hr}\tbr{w}olan	&	{\hr\hr}\tbr{w}ii	&	{\hr\hr}\tbr{w}aʔ	&	{\hr\hr}\tbr{w}eʊ	&	{\hr\hr}\tbr{w}uun	&	\\
Karey	&	{\hr\hr}\tbr{ɸ}o\tbr{ɸ}olo-n	&	{\hr\hr}\tbr{ɸ}olan	&	{\hr\hr}\tbr{w}uaw	&		&	{\hr\hr}\tbr{w}ew	&	{\hr\hr}\tbr{w}uwun	&	\\
Batuley	&	{\hr\hr}\tbr{ɸ}oloi	&	{\hr\hr}\tbr{ɸ}ulan	&	{\hr\hr}\tbr{ɸ}ui	&	{\hr\hr}\tbr{ɸ}at\su{Δ}	&	{\hr\hr}\tbr{ɸ}a\tbr{ɸ}a-ŋ	&	{\hr\hr}\tbr{ɸ}u\tbr{ɸ}un	&	{\hr\hr}\tbr{ɸ}ae\tbr{ɸ}\\
Mariri	&	{\hr\hr}\tbr{ɸ}ul\tbr{ɸ}ulu-n	&	{\hr\hr}\tbr{ɸ}ulan	&	{\hr\hr}\tbr{ɸ}ui	&		&	{\hr\hr}\tbr{ɸ}e\tbr{ɸ}	&	{\hr\hr}\tbr{ɸ}u\tbr{ɸ}un	&	\\
Manombai	&		&	{\hr\hr}\tbr{ɸ}ulan	&	{\hr\hr}\tbr{ɸ}uay	&	{\hr\hr}\tbr{ɸ}atu	&	{\hr\hr}\tbr{ɸ}a\tbr{ɸ}a-y	&		&	{\hr\hr}\tbr{ɸ}a\tbr{ɸ}u\\
NE. Tar. (Fb)\su{†}	&		&	{\hr\hr}\tbr{h}ulan	&		&	{\hr\hr}\tbr{h}ɔt	&	{\hr\hr}\tbr{ɸ}e\tbr{ɸ}a	&		&	\\
NE. Tar. (Gm)\su{†}	&		&	{\hr\hr}\tbr{h}olan	&	{\hr\hr}\tbr{h}uʤa	&	{\hr\hr}\tbr{ʍ}at	&	{\hr\hr}\tbr{ʍ}e\tbr{ʍ}a	&	{\hr\hr}\tbr{ʍ}u\tbr{ʍ}in	&	\\
SE. Tar. (Sr)\su{†}	&		&	{\hr\hr}\tbr{ɸ}ôlan	&	{\hr\hr}\tbr{ɸ}uʤɔ	&	{\hr\hr}\tbr{ɸ}at	&	{\hr\hr}\tbr{ɸ}ɛ\tbr{ɸ}ɔ	&	{\hr\hr}\tbr{ɸ}u\tbr{ɸ}in	&	\\
NW. Tar. (Ju)\su{†}	&		&	{\hr\hr}\tbr{h}ulan	&	{\hr\hr}\tbr{h}uɛ	&	{\hr\hr}\tbr{h}ɔt	&	{\hr\hr}\tbr{h}a\tbr{h}e	&		&	{\hr\hr}\tbr{h}ɔ\tbr{p}\\
SW. Tar. (Kk)\su{†}	&	{\hr\hr}\tbr{ɸ}ul-di\su{ǁ}	&	{\hr\hr}\tbr{ɸ}ôlan	&	{\hr\hr}\tbr{ɸ}uay	&	{\hr\hr}\tbr{ɸ}ɔt	&	{\hr\hr}\tbr{ɸ}a\tbr{ɸ}a-y	&	{\hr\hr}\tbr{ɸ}u\tbr{ɸ}an\su{¶}	&	{\hr\hr}\tbr{ɸ}ɔ\tbr{p}\\
		\lspbottomrule
	\end{tabular}
		\begin{tablenotes}\tnsize
			\item[†]	{Northeast Tarangan (Fb) from Fatlabata,
								Northeast Tarangan (Gm) from Gomar Meti,
								Southeast Tarangan (Sr) from Salarem village, Southwest Tarangan
								(Kk) from Kalar-kalar, Northwest Tarangan (Ju) from Juring.}
			\item[‡]	{Reflexes of **pua- ‘fruit’ are \tsc{3sg} forms except Lola
								\it{ɸua-di} which is \tsc{3pl}.}
			\item[Δ]	{Batuley \it{ɸat} is given as ‘dead weight anchor; anchor
								consisting of a blunt heavy weight such as a stone or piece of concrete’ \citep[251]{Daigle2015}.}
			\item[{ǁ}]	{Dobel \it{ɸulu-}
								and Southwest Tarangan \it{ɸul-di} are ‘eyebrows, eye lashes’.}
			\item[¶]	{East Kola \it{ɸʷuɸʷun}
								and Southwest Tarangan \it{ɸuɸan} are ‘roof ridge’.
								Other reflexes mean ‘top, surface’.
								PMP *bubuŋ is reconstructed meaning ‘ridge of the roof; ridge
								of a mountain, peak; deck of a boat’.}
		\end{tablenotes}
	\end{threeparttable}
\end{adjustbox}
\end{table}


\largerpage
\scp{Internal structure}{06-02}
Within the \tsc{Aru} subgroup there is a \tsc{Nuclear Aru} node
which further contains a \tsc{South Aru} node.
The \tsc{Nuclear Aru} node is defined by *s > \it{h} > {\0},
which occurs in all languages except Ujir.
This leaves Ujir as a family level isolate within \tsc{Aru}.
West Kola has *s > \it{h} while other languages have *s > {\0},
probably through intermediate **h.
Although some languages appear to have *s > \it{y},
this is actually due to glide insertion after loss of *s,
or a shift of final *i > \it{y} after loss of intervocalic *s.
Examples of PMP *s are given in \trf{06-13}.

\begin{table}\tcp{\tsc{Aru} *s}{06-13}
\begin{adjustbox}{width=\textwidth}
	\begin{threeparttable}\st{2pt}
	\begin{tabular}{lllllllll}\lsptoprule
local gl.	&	{\hr}snake	&	{\hr}they	&	{\hr}one 	&	{\hr}what	&	{\hr}sago tree	&	{\hr}sea	&	{\hr}bow	&	{\hr}paddle\\
PMP	&	\hr*\tbr{s}awa\su{†}	&	\hr*\tbr{s}i-ida	&	{\hr}?\hr*ə\tbr{s}a	&		&		&	\hr*ta\tbr{s}ik	&	\hr*bu\tbr{s}uR	&	\hr*bəR\tbr{s}ay\\
Proto-Aru	&	**\tbr{s}awa	&	**\tbr{s}ira	&	**\tbr{s}etu	&	**\tbr{s}a\su{‡}	&	**\tbr{s}akaw	&	**ta\tbr{s}i	&	**pu\tbr{s}uR	&	**peR\tbr{s}ay\\ \midrule
Ujir	&	{\hr\hr}\tbr{s}awa	&	{\hr\hr}\tbr{s}ira	&	{\hr\hr}\tbr{s}et	&	{\hr\hr}\tbr{s}a	&	{\hr\hr}\tbr{s}akaw	&	{\hr\hr}te\tbr{s}	&	{\hr\hr}ɸu\tbr{s}i	&	{\hr\hr}ɸe\tbr{s}a\\
Kompane	&	{\hr\hr}yaa	&		&	{\hr\hr}etu	&	{\hr\hr}ya	&	{\hr\hr}akaw	&		&		&	\\
West Kola	&	{\hr\hr}\tbr{h}awa	&		&	{\hr\hr}\tbr{h}ot	&	{\hr\hr}en \tbr{h}a	&	{\hr\hr}\tbr{h}akaw	&		&		&	\\
East Kola	&	{\hr\hr}awa	&	{\hr\hr}iri	&	{\hr\hr}ot	&	{\hr\hr}ye	&	{\hr\hr}akaw	&	{\hr\hr}tai rei	&	{\hr\hr}ɸuuh	&	{\hr\hr}ɸe\tbr{h}a\su{ǁ}\\
Lola	&	{\hr\hr}yawa	&		&	{\hr\hr}ye	&	{\hr\hr}menayá	&	{\hr\hr}yaʔu	&		&		&	\\
Lorang	&	{\hr\hr}yagwa	&		&	{\hr\hr}yetu	&	{\hr\hr}ya	&	{\hr\hr}yaka	&		&		&	\\
Dobel	&	{\hr\hr}yakwa	&	{\hr\hr}ʔiri, yiri	&	{\hr\hr}e‹t›tú	&	{\hr\hr}ya	&	{\hr\hr}yaʔu	&	{\hr\hr}tay	&	{\hr\hr}ɸiyar	&	{\hr\hr}ɸeri\\
Koba	&	{\hr\hr}yáʍa	&		&	{\hr\hr}ekatú	&	{\hr\hr}ya	&	{\hr\hr}yaka	&	{\hr\hr}tay	&		&	\\
Barakai	&	\hr\cg{(ʤel)}	&		&	{\hr\hr}etí	&	\hr\cg{(woy)}	&	{\hr\hr}aga	&		&		&	\\
Karey	&	\hr\cg{(ʤel)}	&		&	{\hr\hr}etú	&	\hr\cg{(oín)}	&	{\hr\hr}agaw	&		&		&	\\
Batuley	&	{\hr\hr}aw	&	{\hr\hr}id	&	{\hr\hr}et	&	\hr\cg{(ula)}	&	{\hr\hr}aga	&	{\hr\hr}tay	&	{\hr\hr}ɸo(w)or	&	{\hr\hr}ɸer\\
Mariri	&		&		&	{\hr\hr}eʤ	&	\hr\cg{(totóo)}	&	{\hr\hr}ágao	&	{\hr\hr}tay	&		&	\\
Manombai	&		&	{\hr\hr}ira	&	{\hr\hr}etu	&	\hr\cg{(akarey)}	&	{\hr\hr}eka	&		&		&	\\
NE. Tarang.	&	{\hr\hr}yɔ	&	{\hr\hr}ira	&	{\hr\hr}ôt	&	{\hr\hr}ʤa	&	{\hr\hr}aka	&		&		&	\\
SE. Tarang.	&	{\hr\hr}iɔ	&	{\hr\hr}dir	&	{\hr\hr}ɛt-na	&	{\hr\hr}ʤa	&	{\hr\hr}agaw	&		&	{\hr\hr}ɸɔyar	&	\\
NW. Tarang.	&	{\hr\hr}yɔ, iɔ	&	{\hr\hr}ir	&	{\hr\hr}ôt	&	{\hr\hr}ya	&	{\hr\hr}aka	&	{\hr\hr}tay\su{Δ}	&	{\hr\hr}ɸɛar	&	{\hr\hr}ɸɛra\\
SW. Tarang.  	&	\hr\cg{(ʤɛla)}	&	{\hr\hr}êra	&	{\hr\hr}ôt	&	\hr\cg{(lebá)}	&	{\hr\hr}okaw	&	{\hr\hr}tay	&	{\hr\hr}ɸɛr	&	{\hr\hr}ɸera\\
		\lspbottomrule
	\end{tabular}
		\begin{tablenotes}\tnsize
			\item[†]	{PMP *sawa is reconstructed with the meaning ‘python’.}
			\item[‡]	{Likely cognates of **sa ‘what’ occur in \tsc{Timor-Babar} languages:
								Proto-Rote-Meto **saa, Helong \it{saa}, Tetun \it{saa} ‘what’ \citep[319]{Edwards2021},
								though these could be reflexes of PMP *sapa.}
			\item[Δ]	{Northwest Tarangan \it{tay} is from Juring village
								and \it{ɸɛra} ‘paddle’ from Rebi village.}
			\item[{ǁ}]	{Medial \it{h} in East Kola \it{ɸeha} ‘paddle’ is probably
								due to regular **R > \it{h},
								not *s > \it{h}.}
		\end{tablenotes}
	\end{threeparttable}
\end{adjustbox}
\end{table}

Most \tsc{Nuclear Aru} languages also share instances of *t >
**s, normally before front vowels.
However, this change is not complete in most languages,
and exceptions are common.
It probably spread between \tsc{Nuclear Aru} languages by diffusion,
or it may have begun at Proto-Nuclear Aru in a drift-like fashion,
with languages carrying it out to different extents.
Examples of *t > \it{t},
**s are given in \trf{06-14}.
%with instances of *t > **s shaded.
Some languages have subsequent **s > {\0}.
Glides are often subsequently inserted before vowel-initial words
that result from initial **s > {\0}.

\begin{table}\tcp{\tsc{Nuclear Aru} *t > \it{t}, \it{s}}{06-14}
\begin{adjustbox}{width=\textwidth}
	\begin{threeparttable}\st{2pt}
	\begin{tabular}{llllllll}\lsptoprule
local gl.	&	nine	&	wound	&	thick	&	three	&	plant (v.)\su{‡}	&	{\hr}we (incl.)	&	see\\
PMP	&		&		&	bamboo	&		&		&	\hr*kita	&	\\
Proto-Aru	&	**\tbb{t}era\su{†}	&	**\tbb{t}umbu	&	**\tbb{t}emal	&	**la\tbb{t}i	&	**-sa\tbb{t}i	&	**\tbb{t}i\tbb{t}a	&	**-\tbb{t}ewa\\ \midrule
Ujir	&	{\hr\hr}\tbb{t}era	&	{\hr\hr}\tbb{t}ubu	&	{\hr\hr}\tbb{t}emal	&	{\hr\hr}la\tbb{t}i	&	{\hr\hr}sa\tbb{t}	&	{\hr\hr}\tbb{t}ita	&	{\hr\hr}-\tbb{t}o\\
Kompane	&	{\hr\hr}\tbb{t}eri	&	{\hr\hr}\tbr{s}ubu	&		&	{\hr\hr}la\tbr{s}i	&		&		&	{\hr\hr}-\tbr{s}oo\\
East Kola	&	{\hr\hr}\tbb{t}era	&	{\hr\hr}\tbr{s}um	&		&	{\hr\hr}la\tbr{s}	&	{\hr\hr}-aa\tbr{s}	&	{\hr\hr}\tbr{s}ita	&	{\hr\hr}-\tbr{s}oo\\
Lola	&	{\hr\hr}\tbr{s}eri	&		&	{\hr\hr}\tbr{s}eman	&	{\hr\hr}la\tbr{s}i	&		&	{\hr\hr}\tbr{h}ita	&	{\hr\hr}-\tbr{s}ewa\\
Lorang	&	{\hr\hr}\tbr{s}era	&	{\hr\hr}ubu	&	{\hr\hr}yemal	&	{\hr\hr}lay\su{Δ}	&	{\hr\hr}-yay	&	{\hr\hr}ʔita	&	\hr\cg{(m-ɸúlar)}\\
Dobel	&	{\hr\hr}yera	&		&	{\hr\hr}yamal	&	{\hr\hr}lay	&		&	{\hr\hr}ʔita	&	\hr\cg{(-yara)}\\
Koba	&	{\hr\hr}\tbr{s}írey	&		&		&	{\hr\hr}lay\su{Δ}	&		&	{\hr\hr}ʔita	&	\\
Barakai	&	{\hr\hr}\tbr{s}er	&	{\hr\hr}\tbr{s}om	&	{\hr\hr}\tbr{s}emal	&	{\hr\hr}laʔ	&		&	{\hr\hr}\tbr{s}e	&	{\hr\hr}-ʤe\\
Karey	&	{\hr\hr}\tbr{s}er	&		&		&	{\hr\hr}lado	&		&	{\hr\hr}\tbr{s}edo	&	\hr\cg{(-wólar)}\\
Batuley	&	{\hr\hr}\tbr{s}er	&	{\hr\hr}\tbr{s}um	&		&	{\hr\hr}lae\tbr{s}\su{Δ}	&	{\hr\hr}-e\tbr{s}	&	{\hr\hr}\tbr{s}it	&	\hr\cg{(it)}\\
Mariri	&	{\hr\hr}\tbr{s}er(i)	&	{\hr\hr}\tbr{s}um	&		&	{\hr\hr}laʤ	&		&	{\hr\hr}\tbr{s}in	&	\\
Manombai	&	{\hr\hr}\tbb{t}era	&	{\hr\hr}\tbr{s}ubu	&	{\hr\hr}\tbb{t}emal	&	{\hr\hr}la\tbr{s}i	&		&	{\hr\hr}\tbr{s}ita	&	\hr\cg{(-erei)}\\
NE. Tarang.	&	{\hr\hr}\tbr{s}ira	&		&		&	{\hr\hr}la\tbb{t}	&		&	{\hr\hr}\tbr{s}êta	&	\hr\cg{(-ʤala)}\\
SE. Tarang.	&	{\hr\hr}\tbr{s}ɛr	&		&		&	{\hr\hr}la\tbb{t}	&		&	{\hr\hr}\tbr{s}ɛr	&	\hr\cg{(-ŋnan)}\\
NW. Tarang.  	&	{\hr\hr}\tbr{s}êra	&		&	{\hr\hr}\tbb{t}emal	&	{\hr\hr}la\tbb{t}	&		&	{\hr\hr}\tbr{s}ita	&	\hr\cg{(-ʤɔla)}\\
SW. Tarang.  	&	{\hr\hr}\tbr{s}êra	&	{\hr\hr}\tbr{s}um	&		&	{\hr\hr}la\tbb{t}	&	{\hr\hr}-e\tbb{t}	&	{\hr\hr}\tbr{s}êta	&	{\hr\hr}-ʤɔw\\
		\lspbottomrule
	\end{tabular}
		\begin{tablenotes}\tnsize
			\item[†]	{Possible cognates of **tera ‘nine’ occur in Flores; e.g.
								Ngadha \it{təra əsa} ‘nine’.}
			\item[‡]	{The full range of meanings for Proto-Aru **-sati is ‘plant
								with dibble, throw, poke/stab’.}
			\item[Δ]	{Lorang \it{lay},
								Koba \it{lay}, and Batuley \it{laes} are \tsc{3pl} forms.
								Other forms are inanimate.}
		\end{tablenotes}
	\end{threeparttable}
\end{adjustbox}
\end{table}

\citet[131-133]{Collins1982c} identifies a \tsc{South Aru} node,
which excludes Kompane and \tsc{Kola}.
There are two exclusively shared innovations which define \tsc{South Aru}.
The first is merger of Proto-Aru **R/**r > \it{r} (see \trf{06-04} and \trf{06-05}).
The second is *p > **ɸ > **β (see \trf{06-11}),
Additional support for \tsc{South Aru} comes from fortition of **w.
This fortition affects most \tsc{South Aru} words,
though there are sporadic exceptions.
Fortition of **w also occurs in Kompane,
and is thus not exclusive to \tsc{South Aru}.
The result of this fortition is usually **w > \it{gw}.
\tsc{West Tarangan} and Manombai have subsequent de-labialisation
**gw > \it{g}, which is also attested in most Karey words
and some Lorang words.
Dobel has **w > **gw > \it{kw}
and Koba has **w > **gw > \it{ʍ},
which was almost certainly from intermediate **kw.
Thus, Dobel and Koba share *w > **gw > **\it{kw}.
Examples of fortition of **w are given in \trf{06-15}.
%with \tsc{South Aru} exceptions indicated.

\begin{table}
\caption[Proto-Aru **w]{Proto-Aru **w\su{†}}\label{tab:06-15}
\begin{adjustbox}{width=\textwidth}
	\begin{threeparttable}\st{2pt}
	\begin{tabular}{ll@{ }l@{ }l@{ }l@{ }l@{ }l@{ }l}\lsptoprule
local gl.	&	{\hr}water	&	{\hr}root	&	{\hr}ySib	&	dry in sun	&	{\hr}rattan	&	wake up	&	sleepy\\
PMP	&	\hr*\tbr{w}ahiyəR	&	\hr*\tbr{w}akaR	&	\hr*huaji	&	\hp{*-}*\tbr{w}aRi	&		&		&	\\
P.-Aru	&	**\tbr{w}ayaR	&	**\tbr{w}akaR	&	**\tbr{w}ali	&	**-\tbr{w}aRay	&	**\tbr{w}aŋguR	&	**-\tbr{w}aŋaR	&	**-\tbr{w}oŋga\\ \midrule
Ujir	&	{\hr\hr}\tbr{w}ay	&	{\hr\hr}\tbr{w}aki	&	{\hr\hr}\tbr{w}ɪl	&	{\hr\hr}-\tbr{w}oy	&	{\hr\hr}\tbr{w}aŋi	&	{\hr\hr}-\tbr{w}aŋay	&	\\
Kompane	&	{\hr\hr}\tbr{gw}ah	&	{\hr\hr}\tbr{w}aki	&	{\hr\hr}\tbr{gw}áel	&	{\hr\hr}-\tbr{gw}ahi	&	{\hr\hr}\tbr{gw}aŋu	&	{\hr\hr}-\tbr{gw}aŋah	&	{\hr\hr}-\tbr{g}oŋa\\
East Kola	&	{\hr\hr}\tbr{w}eeh	&	{\hr\hr}\tbr{w}akih	&	{\hr\hr}\tbr{w}el	&	{\hr\hr}-\tbr{w}eh	&	{\hr\hr}\tbr{w}aguh	&	{\hr\hr}-\tbr{w}aŋah	&	{\hr\hr}-\tbr{w}oga\\
Lola	&	{\hr\hr}\tbr{w}aʔ	&	{\hr\hr}\tbr{w}aʔi	&	{\hr\hr}\tbr{w}en	&	{\hr\hr}-\tbr{w}ahi	&	{\hr\hr}\tbr{w}aŋ	&	{\hr\hr}-\tbr{w}aŋ	&	\\
Lorang	&	{\hr\hr}\tbr{gw}ar	&	{\hr\hr}\tbr{gw}akir	&	{\hr\hr}\tbr{g}el	&	{\hr\hr}-\tbr{g}era	&		&	{\hr\hr}-\tbr{gw}aŋar	&	{\hr\hr}-\tbr{g}oŋa\\
Dobel	&	{\hr\hr}\tbr{kw}ar	&	{\hr\hr}\tbr{kw}aʔir	&	{\hr\hr}\tbr{kw}ali-	&	{\hr\hr}-\tbr{kw}ari	&	{\hr\hr}\tbr{kw}aŋur	&	{\hr\hr}-\tbr{kw}aŋar	&	\\
Koba	&	{\hr\hr}\tbr{ʍ}ar	&	{\hr\hr}\tbr{ʍ}akir	&	{\hr\hr}\tbr{ʍ}éla	&	{\hr\hr}-\tbr{ʍ}era	&	{\hr\hr}\tbr{ʍ}aŋar	&	{\hr\hr}-\tbr{ʍ}áŋar	&	\\
Barakai	&	{\hr\hr}\tbr{gw}aor	&	{\hr\hr}\tbr{gw}áer	&	{\hr\hr}\tbr{gw}el-	&	\hp{*-}\cg{(da-bálar)}	&	{\hr\hr}\tbr{gw}agur\su{Δ}	&		&	{\hr\hr}-\tbr{gw}in\\
Karey	&	{\hr\hr}\tbr{g}aor	&	{\hr\hr}\tbr{g}á\tbr{g}er	&	{\hr\hr}\tbr{g}óliŋ	&	{\hr\hr}-\tbr{ɸ}er	{\cgr}&	{\hr\hr}\tbr{gw}aŋur	&		&	\\
Batuley	&	{\hr\hr}\tbr{gw}ayor	&	{\hr\hr}\tbr{gw}ágwer	&	{\hr\hr}\tbr{gw}al	&	{\hr\hr}\tbr{w}ayer	{\cgr}&	{\hr\hr}\tbr{gw}aŋor	&	{\hr\hr}-\tbr{w}aŋar	{\cgr}&	\\
Mariri	&	{\hr\hr}\tbr{gw}ayor	&	{\hr\hr}\tbr{w}ager	{\cgr}&	{\hr\hr}\tbr{gw}el-	&	{\hr\hr}-\tbr{w}er	&	{\hr\hr}\tbr{gw}aŋor	&		&	\\
Manombai	&	{\hr\hr}\tbr{g}ayar	&	{\hr\hr}\tbr{g}ákar-i	&	{\hr\hr}\tbr{g}ali	&	{\hr\hr}-\tbr{g}aray	&		&	{\hr\hr}-\tbr{g}aŋar	&	{\hr\hr}-\tbr{g}oŋa\\
NE. Tar.	&	{\hr\hr}\tbr{gw}ar	&	{\hr\hr}\tbr{gw}akir	&	{\hr\hr}\tbr{gw}ɛl-	&	{\hr\hr}-\tbr{gw}arɛ	&	{\hr\hr}\tbr{gw}aŋar	&		&	{\hr\hr}-\tbr{g}ɔm\\
SE. Tar.	&	{\hr\hr}\tbr{gw}ar	&	{\hr\hr}\tbr{gw}agwir	&	{\hr\hr}\tbr{gw}ɛl-	&	{\hr\hr}-\tbr{gw}aray	&	{\hr\hr}\tbr{gw}aŋar	&		&	{\hr\hr}-\tbr{g}ɔm{\tl}gɔm\\
NW. Tar.	&	{\hr\hr}\tbr{g}ar	&	{\hr\hr}\tbr{g}akir	&	{\hr\hr}ʤɛl-\su{‡}	&	{\hr\hr}-\tbr{g}aray	&	{\hr\hr}\tbr{g}aŋar	&		&	{\hr\hr}-\tbr{g}ɔm\\
SW. Tar.	&	{\hr\hr}\tbr{g}ar	&	{\hr\hr}\tbr{g}akar	&	{\hr\hr}ʤɛl-	&	{\hr\hr}-\tbr{g}aray	&	{\hr\hr}\tbr{g}aŋar	&	{\hr\hr}-\tbr{g}aŋar	&	{\hr\hr}-\tbr{g}ɔm\\
		\lspbottomrule
	\end{tabular}
		\begin{tablenotes}\tnsize
			\item[†]	{Shading indicates \tsc{South Aru} words which do not have fortition of *w.}
			\item[‡]	{\tsc{West Tarangan} \it{ʤɛl-} ‘younger sibling’ has palatalisation
								of **g before a front vowel.}
			\item[Δ]	{Barakai \it{gwagur} is from Mesiang village.}
		\end{tablenotes}
	\end{threeparttable}
\end{adjustbox}
\end{table}

\citet{Nivens2017} includes Lola within \tsc{South Aru},
but Lola does not have fortition of *w
and does not merge **R/**r > \it{r}.
Instead, Lola retains **r = \it{r}
and has **R > \it{h}.
The change **R > \it{h} also occurs in Kompane
and \tsc{Kola}, and on this basis,
we could place these three languages in a single node.
However, West Kola has *s > \it{h} (see \trf{06-13}), while East
Kola, Kompane, and Lola have *s > (**h) > {\0}.
Thus, if we posited **R > \it{h} at a single node encompassing
these languages, it would incorrectly predict a merger of *j/*R/*s > **R/**h >
{\0} in these languages.
Instead, it seems that the change **R > \it{h} in Lola is independent
of that in \tsc{Kola} and Kompane.
\tsc{Kola} and Kompane are spoken adjacent to one another
and we can posit that **R > \it{h} spread among these languages
by diffusion after East Kola
and Kompane had undergone (*s >) **h > {\0}.
Because no other significant sound changes are shared between
these languages, we place \tsc{Kola}, Kompane,
and Lola as isolates within \tsc{Nuclear Aru}.

Innovation-defined nodes within \tsc{South Aru} are not easy to identify.
Lorang, Dobel, and Koba all share another round of word-initial
glide epenthesis (after the fortition of glides inserted at an
earlier stage), and could be placed together on the basis of this change.
However, this change is also found in Lola which we exclude from \tsc{South Aru}.
Furthermore, glide epenthesis is a recurrent sound change in
the languages of Aru, which has probably been ongoing ever since
the formation of this subgroup (see \srf{06-01-03}).

We identify only three further innovation-defined nodes within \tsc{South Aru}.
Firstly, Dobel and Koba clearly group together on the basis of
*y/*z > **y > **ʤ > \it{s} (see examples in \trf{06-06},
and \trf{06-08}), as well as devoicing of *w > **gw.
Dobel has **gw > \it{kw},
and Koba has **w > **gw > \it{ʍ} (voiceless labio-velar approximant).
This group also matches the lexicostatistical classification
of \citet{Hughes1987} in which Koba was lexically closest to Dobel.

Secondly, Barakai and Karey can be grouped together on the basis
of *b > **p > \it{w} in most words (see \trf{06-11}).
This change is different from most other \tsc{Aru} languages
which have *b > \it{ɸ.} Note,
however, that Karey has *p > \it{ɸ} /{\gap}o.
Again, this classification matches the lexical classification of \citet{Hughes1987}.

Thirdly, we place the \tsc{Tarangan} languages together on the
basis of initial PMP *pa > Proto-Aru **ɸ > Proto-South Aru **βa
> \it{ɔ} in most words.
Two examples of *pa > \it{ɔ} in \tsc{Tarangan} have been given
in \trf{06-11} (*paniki > East Tarangan \it{ɔnít} ‘flying fox’
and *pana > \tsc{West Tarangan} \it{-ɔn} ‘shoot’).
Two additional examples are *panaw ‘whitish skin disorder’ >
Southwest Tarangan \it{ɔnaw} (East Kola \it{ɸʷanaw},
Dobel \it{wanu}), and *paŋdan > Southwest Tarangan \it{ɔdan} ‘pandanus’.
Supporting evidence for \tsc{Tarangan} may come from
the innovation of seven-vowel systems in which mid-high ‹ê›
/e/ and ‹ô› /o/ contrast with mid-low /ɛ/ and /ɔ/.

No further innovation-defined nodes are well supported within
\tsc{South Aru} on the basis of the phonological history.
\citet{Hughes1987} treated Mariri as a dialect of Batuley,
but we find little support for this from the historical phonology.
They do share medial *-k- > \it{-g-},
but this change is also found in Karey.
Other changes that are shared by Batuley
and Mariri, such as *p > {\0},
are common among \tsc{Aru} languages.

One unusual change which deserves mention occurs in Mariri,
which has *n > \it{d} word initially and medially.
Examples include: *anak > \it{adeʤ} ‘child’,
*naŋi > \it{-den} ‘swim’,
*nipən > \it{deen} ‘tooth’,
*bunuq > \it{-ɸʊdʊn} ‘kill’,
*kunij > \it{kódel} ‘yellow’,
*niuR > \it{dor} ‘coconut’,
*panua > \it{ɸadú} ‘village’,
and *haŋin > **anin > \it{aden} ‘wind’.


\scp{Final vowel-glide sequences}{06-03}
The reflexes of final vowel-glide sequences *-aw
and *-ay in the Aru languages are diverse.
 In this section we only focus on the reflexes of final vowel-glide
sequences for words in which a medial consonant was preserved,
as the reflexes are clearest in this environment.
Reflexes of vowel-glide sequences include: preservation of the
full sequence (*-aw = \it{aw;} *-ay = \it{ay}),
fusion to a mid vowel (*-ay > \it{ɛ}),
deletion of the final glide (*-aw \it{>} -\it{a}),
and preservation of the final glide as a vowel (*-aw > \it{u}; *-ay > \it{i}).
The full range of reflexes found throughout the Wallacea region
are seen in different \tsc{Aru} languages.
No changes to vowel-glide sequences align with different nodes
within \tsc{Aru} and a single language does not necessarily have parallel changes
for *-aw and *-ay.
Examples of the reflexes of *-aw are given in \trf{06-16},
and *-ay in \trf{06-17}.

There is not enough data to determine a consistent reflex of
*-aw for Lorang, Barakai, Karey, Mariri, and Manombai.
There is also not enough data to determine a consistent reflex
of *-ay for Kompane,
Lola, Lorang, Koba, Barakai,
Karey, Mariri, and East Tarangan.
(Additionally, the reflexes of *bəRsay ‘paddle’ (\trf{06-13})
appear to be irregular with respect to the final vowel-glide
sequence.) 

\begin{table}\tcp{\tsc{Aru} *-aw = **-aw}{06-16}
\begin{adjustbox}{width=\textwidth}
	\begin{threeparttable}\st{2pt}
	\begin{tabular}{lllllll}\lsptoprule
local gl.	&		&	sago	&	\hp{*-}walk	&	{\hr}sun	&	\hp{*-}steal	&	house\\
PMP	&	\hr*-aw	&		&	\hp{*-}*lak\tbr{aw}	&	\hr*qaləj\tbr{aw}	&	\hp{*-}*nak\tbr{aw}	&	\\
Proto-Aru	&	**-aw	&	**sak\tbr{aw}	&	**-lak\tbr{aw}	&	**laR\tbr{aw}	&	**-naŋg\tbr{aw}	&	**pal\tbr{aw}\\ \midrule
Ujir	&	{\hr\hr}-aw, -u	&	{\hr\hr}sak\tbr{aw}	&		&	{\hr\hr}law\tbr{u}	&	{\hr\hr}-y-naŋ\tbr{aw}	&	\\
Kompane	&	{\hr\hr}-aw	&	{\hr\hr}ak\tbr{aw}	&		&	{\hr\hr}lah\tbr{aw}	&		&	\\
East Kola	&	{\hr\hr}-aw	&	{\hr\hr}ak\tbr{aw}	&	{\hr\hr}-lak\tbr{áu}	&	{\hr\hr}lah\tbr{aw}	&	{\hr\hr}-fanag\tbr{aw}	&	{\hr\hr}ɸal\tbr{aw}\\
Lola	&	{\hr\hr}-u , -a	&	{\hr\hr}yaʔ\tbr{u}	&		&	{\hr\hr}lah\tbr{u}	&		&	{\hr\hr}wal\tbr{a}sín\\
Lorang	&	{\hr\hr}(-a)	&	{\hr\hr}yak\tbr{a}	&		&		&		&	\\
Dobel	&	{\hr\hr}-u	&		&		&	{\hr\hr}lar\tbr{u}\su{‡}	&	{\hr\hr}-naŋ\tbr{u}	&	{\hr\hr}ɸal\tbr{u}\su{Δ}\\
Koba	&	{\hr\hr}-a	&	{\hr\hr}yak\tbr{a}	&		&	{\hr\hr}lar\tbr{a}	&		&	\\
Barakai	&	{\hr\hr}(-a)	&	{\hr\hr}ag\tbr{a}	&		&		&		&	\\
Karey	&	{\hr\hr}(-aw)	&	{\hr\hr}ag\tbr{aw}	&		&		&		&	\\
Batuley	&	{\hr\hr}-a	&	{\hr\hr}ag\tbr{a}	&		&	{\hr\hr}taɸer lar\tbr{a}	&		&	\\
Mariri	&	{\hr\hr}(-ao)	&	{\hr\hr}ág\tbr{ao}	&		&		&		&	\\
Manombai	&	{\hr\hr}(-a)	&	{\hr\hr}ek\tbr{a}	&		&		&		&	\\
NE. Tar.\su{†}	&	{\hr\hr}-a	&	{\hr\hr}ak\tbr{a}	&	{\hr\hr}-r-lɔk\tbr{a}	&		&	{\hr\hr}-naŋ\tbr{a}	&	{\hr\hr}hál\tbr{a}\\
SE. Tar.	&	{\hr\hr}-aw	&	{\hr\hr}ag\tbr{aw}	&	{\hr\hr}-r-lag\tbr{aw}	&		&		&	{\hr\hr}ɸal\tbr{aw}\\
NW. Tar. 1	&	{\hr\hr}-a	&	{\hr\hr}ak\tbr{a}	&	{\hr\hr}-r-lɔk\tbr{a}	&		&	{\hr\hr}-nɔŋ\tbr{a}	&	\\
NW. Tar. 2	&	{\hr\hr}-a	&	{\hr\hr}ak\tbr{a}	&	{\hr\hr}-r-lɔk\tbr{a}	&		&	{\hr\hr}-nɔŋ\tbr{a}	&	\\
SW. Tar. 1	&	{\hr\hr}-aw	&	{\hr\hr}ɔk\tbr{aw}	&	{\hr\hr}-r-lɔk\tbr{aw}	&		&	{\hr\hr}-nɔŋ\tbr{aw}	&	{\hr\hr}ɸal{\tl}ɸal\tbr{aw}\su{ǁ}\\
SW. Tar. 2	&	{\hr\hr}-aw	&	{\hr\hr}ɔk\tbr{aw}	&		&		&	{\hr\hr}-nɔŋ\tbr{aw}	&	\\
		\lspbottomrule
	\end{tabular}
		\begin{tablenotes}\tnsize
			\item[†]	{Data from varieties of \tsc{Tarangan}
								are as follows: NE. Tarangan, Jelia village;
								SE. Tarangan from Dosi Mar village;
								NW. Tarangan 1 from Doka Timur village;
								NW. Tarangan 2 from Lutur village;
								SW. Tarangan 1 from Ngaibor village; and
								SW. Tarangan 2 from Kalar-kalar village.}
			\item[‡]	{Lola \it{lahu} ‘sun’ is from Warabal village.
								Lola village has \it{laɸo} ‘sun’.}
			\item[Δ]	{Dobel \it{ɸalu} is ‘very large clan house’.}
			\item[{ǁ}]	{Rick Nivens (pers. comm.
								April 2024) doubts that Ngaibor \it{ɸal{\tl}ɸalaw} is a valid form
								and proposes that it is probably due to contact with SE. Tarangan.}
		\end{tablenotes}
	\end{threeparttable}
\end{adjustbox}
\end{table}

\begin{table}\tcp{\tsc{Aru} *-ay = **-ay}{06-17}
\begin{adjustbox}{width=\textwidth}
	\begin{threeparttable}\st{2pt}
	\begin{tabular}{llllllll}\lsptoprule
local gl.	&		&	{\hr}dry in sun	&	{\hr}climb	&	bitter	&	open	&	boat	&	sailboat\\
PMP	&	\hr*-ay	&	\hr*waRi	&	\hr*sak\tbr{ay}	&		&		&		&	\\
Proto-Aru	&	**-ay	&	**-waR\tbr{ay}	&	**sak\tbr{ay}	&	**mak\tbr{ay}	&	**pak\tbr{ay}	&	**kal\tbr{ay}	&	**let\tbr{ay}\\	\midrule
Ujir	&	{\hr\hr}-i, -ay	&	{\hr\hr}-wo\tbr{y}	&		&		&	{\hr\hr}ɸak\tbr{i}	&	{\hr\hr}kal\tbr{i}	&	{\hr\hr}let\tbr{ay}\\
Kompane	&	{\hr\hr}(-i)	&	{\hr\hr}-gwah\tbr{i}	&		&		&		&		&	\\
East Kola	&	{\hr\hr}-i	&	{\hr\hr}-weh	&	{\hr\hr}-as\tbr{i}	&	{\hr\hr}mak\tbr{i}	&	{\hr\hr}-ɸak\tbr{i}	&	{\hr\hr}kal\tbr{i}	&	{\hr\hr}let\tbr{i}\\
Lola	&	{\hr\hr}(-i)	&	{\hr\hr}-wah\tbr{i}	&		&		&		&		&	\\
Lorang	&	{\hr\hr}(-a)	&	{\hr\hr}-ger\tbr{a}	&		&		&		&		&	\\
Dobel	&	{\hr\hr}-ay	&	{\hr\hr}-kwar\tbr{i}	&	{\hr\hr}s{\tl}saʔ\tbr{ay}	&	{\hr\hr}maʔ\tbr{ay}	&	{\hr\hr}-ɸaʔ\tbr{ay}	&	{\hr\hr}ʔal\tbr{ay}	&	{\hr\hr}let\tbr{ay}\\
Koba	&	{\hr\hr}(-a)	&	{\hr\hr}-ʍer\tbr{a}	&		&		&		&		&	\\
Batuley	&	{\hr\hr}-ey	&	{\hr\hr}wa\tbr{ye}r\su{‡}	&		&	{\hr\hr}mag\tbr{ey}	&		&	{\hr\hr}kál\tbr{ey}	&	{\hr\hr}let\tbr{ey}\\
Manombai	&	{\hr\hr}-ay	&	{\hr\hr}-gar\tbr{ay}	&		&		&	{\hr\hr}ɸak\tbr{ay}	&		&	\\
NE. Tar.\su{†}	&	{\hr\hr}(-ɛ)	&	{\hr\hr}-gwar\tbr{ɛ}	&		&		&		&		&	\\
SE. Tar.	&	{\hr\hr}(-ay)	&	{\hr\hr}-gwar\tbr{ay}	&		&		&		&		&	\\
NW. Tar. 1	&	{\hr\hr}-ɛ	&	{\hr\hr}-gar\tbr{ɛ}	&		&	{\hr\hr}mak\tbr{ɛ}	&	{\hr\hr}-ɸák\tbr{ɛ}	&		&	{\hr\hr}lɛt\tbr{ɛ}\\
NW. Tar. 2	&	{\hr\hr}-ay	&	{\hr\hr}-gar\tbr{ay}	&	{\hr\hr}-ak\tbr{ay}	&		&		&		&	{\hr\hr}lɛt\tbr{ay}\\
SW. Tar. 1	&	{\hr\hr}-ay	&	{\hr\hr}-gar\tbr{ay}	&	{\hr\hr}-ak\tbr{ay}	&	{\hr\hr}mak\tbr{ay}	&		&		&	{\hr\hr}lɛt\tbr{ay}\\
SW. Tar. 2	&	{\hr\hr}-ay	&	{\hr\hr}-gar\tbr{ay}	&	{\hr\hr}-ak\tbr{ay}	&	{\hr\hr}mak\tbr{ay}	&	{\hr\hr}-ɸak\tbr{ay}	&	{\hr\hr}kal\tbr{ay}\su{Δ}	&	{\hr\hr}lɛt\tbr{ay}\\
		\lspbottomrule
	\end{tabular}
		\begin{tablenotes}\tnsize
			\item[†]	{\tsc{Tarangan} data is from the same varieties
								and villages as given in \trf{06-16}.
								Northwest Tarangan \it{lɛtay} ‘sailboat’ is from Rebi village.}
			\item[‡]	{Batuley \it{wayer} ‘dry in sun’ has metathesis
								of the medial and final consonants from **warey.}
			\item[Δ]	{Southwest Tarangan has \it{kalin} ‘boats’ as a plural.
								Rick Nivens (pers. comm.
								March 2021) further points out that this is unusual since the
								plural suffix \it{-in} is usually only found on inalienably possessed nouns.}
		\end{tablenotes}
	\end{threeparttable}
\end{adjustbox}
\end{table}

\clearpage\newpage
\scp{\tsc{Aru} classification}{06-04}
All languages have the following sound changes unless otherwise
noted: *q/*h > {\0}; *k = \it{k}; *b > **p > \it{ɸ}; *d > **r
= \it{r}; *z/*y > **y > \it{ʤ} /{\#}{\gap}; *m = \it{m}; *n =
\it{n}; and *l = \it{l}.
The reflexes of *z/*y > **y intervocalically (\srf{06-01-02}),
final vowel-glide sequences (\srf{06-03}),
and Proto-South Aru **β (\srf{06-01-04}) involve multiple splits
and are not usually given here to avoid clutter.
\mbox{}\\\noindent
\dirtree{%This  is needed
.1 {\tsc{Aru} *j/*R > **R; *z/*y/*iV > **y; *d > **r; *p > **ɸ;
							*b > **p; {\0} > \it{y, w,} {\0} /{\#\gap}V}.
		.2 {Ujir [udj] **R > \it{r, ʤ} /{\#\gap}; **R > \it{y,} {\0} /V{\gap}V, /{\gap\#}}.
		.2 {\tsc{Nuclear Aru} *s > **h (only retained as \it{h} in West Kola); *t > \it{t, s}}.
				.3 {\tsc{Kola} [kvv] **ɸ > \it{f}; **p > \it{ɸʷ}}.
						.4 {West Kola **R > \it{h}; **y > \it{ʤ}; (**h = \it{h})}.
						.4 {East Kola **R > \it{h}; **h > {\0}; (**y = \it{y} /{\#\gap})}.
				.3 {Kompane [kvp] **R > \it{h}; **ɸ > {\0} /{\#}{\gap}; **ɸ > \it{w}, {\0} /V{\gap}V}.
				.3 {Lola [lcd] **R > \it{h}; *l > \it{n} /{\gap}{\#}; **ɸ > h /{\#}{\gap};
											 **ɸ > {\0}, \tsc{+round}, \it{w} /V{\gap}V}.
				.3 {\tsc{South Aru} (*p >) **ɸ > **β; **R/**r > \it{r}; *w > **gw; **p > \it{ɸ}}.
						.4 {Manombai [woo] **gw > \it{g}}.
						.4 {Batuley [bay] *k > \it{g} /V{\gap}V}.
						.4 {Mariri [mqi] *n > \it{d} /{\gap}ˈV; *k > \it{g} /V{\gap}V}.
						.4 {\tsc{Dobel-Koba} (**y >) **ʤ > \it{s}; **gw > **kw}.
								.5 {Dobel [kvo] *k > \it{ʔ} (**kw = \it{kw})}.
								.5 {Koba [kpd] **kw > \it{ʍ}}.
						.4 {Lorang [lrn] **gw > \it{gw, g}}.
						.4 {\tsc{Barakai-Karey}}.
								.5 {Barakai [baj]}.
								.5 {Karey [kyd] **gw > \it{g}; *k > \it{g} /V{\gap}V}.
						.4 {\tsc{Tarangan} (*pa >) **βa > \it{ɔ} /{\#}{\gap}; (*b >)
																**p > \it{ɸ {\tl} ʍ {\tl} h}
																(variation \hp{**}between varieties); **p = \it{p} /{\gap}{\#}}.
								.5 {\tsc{West Tarangan} [txn] **gw > \it{g}}.
										.6 {Northwest Tarangan}.
										.6 {Southwest Tarangan}.
								.5 {\tsc{East Tarangan} [tre]}.
										.6 {Northeast Tarangan *k > \it{g} /V{\gap}V}.
										.6 {Southeast Tarangan}.
}%Ensure there is no carriage return before this }

